\documentclass[12pt,openany]{book}
\usepackage[utf8]{inputenc}
\usepackage[T1]{fontenc}
\usepackage[spanish,es-noshorthands]{babel}

\usepackage{amsmath, amssymb, amsthm}
\usepackage{graphicx}
\usepackage{listings}
\usepackage{xcolor}
\usepackage{algorithm}
\usepackage{algpseudocode}
\usepackage{hyperref}
\usepackage{enumitem}
\usepackage{multirow}
\usepackage{array}
\usepackage{booktabs}
\usepackage{caption}
\usepackage{subcaption}
\usepackage{pgfplots}
\usepackage{tikz}
\usepackage{float}

\usetikzlibrary{babel, arrows.meta, positioning}
\pgfplotsset{compat=1.18}

% Listings para Haskell
\lstset{
  language=Haskell,
  basicstyle=\small\ttfamily,
  keywordstyle=\color{blue},
  commentstyle=\color{green!50!black},
  stringstyle=\color{red!70!black},
  numbers=left,
  numberstyle=\tiny\color{gray},
  stepnumber=1,
  breaklines=true,
  frame=single,
  tabsize=2,
  showspaces=false,
  showstringspaces=false
}

% Listings para Python
\lstdefinestyle{python}{
  language=Python,
  basicstyle=\small\ttfamily,
  keywordstyle=\color{blue}\bfseries,
  commentstyle=\color{green!50!black}\itshape,
  stringstyle=\color{red!70!black},
  numberstyle=\tiny\color{gray},
  numbers=left,
  stepnumber=1,
  backgroundcolor=\color{gray!5},
  frame=single,
  breaklines=true,
  tabsize=4
}

% Algoritmos en español
\algrenewcommand\algorithmicrequire{\textbf{Entrada:}}
\algrenewcommand\algorithmicensure{\textbf{Salida:}}
\floatname{algorithm}{Algoritmo}

% Teoremas
\theoremstyle{definition}
\newtheorem{definicion}{Definición}[chapter]
\newtheorem{ejemplo}[definicion]{Ejemplo}
\newtheorem{observacion}[definicion]{Observación}
\newtheorem{ejercicio}[definicion]{Ejercicio}

\theoremstyle{plain}
\newtheorem{teorema}[definicion]{Teorema}
\newtheorem{corolario}[definicion]{Corolario}
\newtheorem{lema}[definicion]{Lema}
\newtheorem{proposicion}[definicion]{Proposición}

% Comandos útiles
\newcommand{\N}{\mathbb{N}}
\newcommand{\Z}{\mathbb{Z}}
\newcommand{\Q}{\mathbb{Q}}
\newcommand{\R}{\mathbb{R}}
\newcommand{\C}{\mathbb{C}}
\newcommand{\Pbb}{\mathcal{P}}

\title{Estructuras Discretas II}
\author{Dr(c) Wilber Ramos Lovón}
\date{2026}

\begin{document}

\frontmatter
\maketitle
\chapter*{Prólogo}
\markboth{Prólogo}{Prólogo}

\section*{¿Por qué este libro?}
Cuando un estudiante de Ciencia de la Computación se enfrenta por primera vez a las matemáticas discretas, suele hacer dos preguntas fundamentales:
\begin{itemize}
  \item ¿para qué sirve esto?
  \item ¿cómo se relaciona con la programación que ya conozco?
\end{itemize}
Este libro es mi intento de responder a esas preguntas de la manera más clara y práctica posible.

\section*{Estructura del libro}
Este libro es la continuación natural de \emph{Estructuras Discretas I} y está organizado en cinco partes principales:
\begin{enumerate}
  \item \textbf{Combinatoria}: El arte del conteo. Desde los principios básicos hasta las aplicaciones en el análisis de algoritmos.
  \item \textbf{Teoría de Grafos}: El estudio de redes y relaciones. Desde las definiciones fundamentales hasta los algoritmos de optimización.
  \item \textbf{Criptografía}: La ciencia de la comunicación segura. Desde los cifrados clásicos hasta el sistema RSA.
  \item \textbf{Programación Funcional con Haskell}: Una introducción práctica a Haskell como herramienta para implementar y verificar conceptos matemáticos.
  \item \textbf{Guías de Práctica}: Ejercicios y proyectos que integran todos los conceptos aprendidos.
\end{enumerate}

\section*{¿Por qué Haskell?}
Haskell no es solo otro lenguaje de programación. Es una herramienta que permite:
\begin{itemize}
  \item Escribir funciones que son matemáticamente puras.
  \item Razonar sobre el código mediante demostraciones formales.
  \item Implementar estructuras de datos con una elegancia que refleja su definición matemática.
  \item Verificar experimentalmente los conceptos teóricos.
\end{itemize}

\section*{Agradecimientos}
A los estudiantes de la Escuela Profesional de Ciencia de la Computación, que con sus preguntas y entusiasmo han enriquecido este material.

\vspace{1cm}
\begin{flushright}
Wilber Roberto Ramos Lovón\\
Arequipa, 2026
\end{flushright}
\tableofcontents

\mainmatter

\part{Técnicas Avanzadas de Conteo}
\chapter{Combinatoria}

Aunque la enumeración o conteo se percibe como un concepto básico de la aritmética, su importancia y complejidad suelen subestimarse. Es esencial comprender los principios fundamentales en lugar de confiar únicamente en fórmulas.

\section{Principio de la multiplicación}

Si una actividad se puede realizar en $t$ pasos sucesivos, y el paso $i$ se puede hacer de $n_i$ maneras, entonces el número total de actividades es $n_1 \cdot n_2 \cdots n_t$.

\begin{ejemplo}
En una picantería se ofrece: 2 entradas, 3 platos de fondo y 4 bebidas. ¿Cuántas formas hay de seleccionar un plato de fondo y una bebida?
\[
\text{Respuesta} = 3 \cdot 4 = 12.
\]
¿Y una entrada, un plato y una bebida?
\[
\text{Respuesta} = 2 \cdot 3 \cdot 4 = 24.
\]
\end{ejemplo}

\begin{ejemplo}
¿Cuántas cadenas de longitud 4 se pueden formar con las letras ABCDE sin repetición?
\begin{enumerate}[label=\alph*)]
  \item Total: $5 \cdot 4 \cdot 3 \cdot 2 = 120$.
  \item Comienzan con B: $1 \cdot 4 \cdot 3 \cdot 2 = 24$.
  \item No comienzan con B: $120 - 24 = 96$.
\end{enumerate}
\end{ejemplo}

\begin{ejemplo}
¿Cuántos subconjuntos tiene $A = \{a_1, \ldots, a_n\}$?
\[
\underbrace{2 \cdot 2 \cdots 2}_{n \text{ factores}} = 2^n.
\]
\end{ejemplo}

\section{Principio de la Suma}

Si $X_1, \ldots, X_t$ son conjuntos ajenos por pares con $n_i$ elementos cada uno, el número de elementos en la unión es $n_1 + \cdots + n_t$.

\begin{ejemplo}
¿Cuántas cadenas de ocho bits comienzan con 101 o con 111?
\[
2^5 + 2^5 = 64.
\]
\end{ejemplo}

\begin{definicion}
Una \textbf{permutación} de $n$ elementos distintos es un ordenamiento de los $n$ elementos.
\end{definicion}

\begin{teorema}
Existen $n!$ permutaciones de $n$ elementos.
\end{teorema}
\begin{proof}
Por el principio de la multiplicación:
$n(n-1)(n-2)\cdots 3\cdot 2\cdot 1 = n!$.
\end{proof}

\begin{definicion}
Una \textbf{permutación $r$} de $n$ elementos es un ordenamiento de $r$ de ellos. Se denota $P(n,r)$.
\end{definicion}

\begin{teorema}
\[
P(n,r) = \frac{n!}{(n-r)!}, \quad r \leq n.
\]
\end{teorema}

\begin{ejemplo}
¿De cuántas maneras se puede elegir presidente, vicepresidente, secretario y tesorero entre 10 delegados?
\[
P(10,4) = 10\cdot 9\cdot 8\cdot 7 = 5040.
\]
\end{ejemplo}

\begin{definicion}
Una \textbf{combinación $r$} de $X$ es una selección no ordenada de $r$ elementos de $X$. Se denota $C(n,r)$ o $\binom{n}{r}$.
\end{definicion}

\begin{teorema}
\[
C(n,r) = \frac{n!}{(n-r)!\,r!}.
\]
\end{teorema}

\begin{ejemplo}
¿De cuántas maneras elegir 3 representantes entre 5 delegados?
\[
C(5,3) = \frac{5!}{2!\,3!} = 10.
\]
\end{ejemplo}

\begin{ejemplo}
Comité de 2 mujeres y 3 hombres de 5 mujeres y 6 hombres:
\[
C(5,2)\,C(6,3) = 10 \cdot 20 = 200.
\]
\end{ejemplo}

\begin{ejemplo}
¿Cuántas cadenas de 8 bits contienen exactamente 4 unos?
\[
C(8,4) = 70.
\]
\end{ejemplo}

\begin{ejemplo}
Baraja de 52 cartas:
\begin{enumerate}[label=\alph*)]
  \item Manos de póquer: $C(52,5) = 2\,598\,960$.
  \item Todas del mismo palo: $4 \cdot C(13,5) = 5148$.
  \item Tres de una denominación y dos de otra: $13 \cdot 12 \cdot C(4,3) \cdot C(4,2) = 3744$.
\end{enumerate}
\end{ejemplo}
\chapter{Relaciones de Recurrencia}

\section{Definiciones}

\begin{definicion}
Una \textbf{relación de recurrencia} para la sucesión $a_0, a_1, a_2, \ldots$ es una ecuación que relaciona $a_n$ con predecesores. Las \textbf{condiciones iniciales} son valores dados explícitamente para un número finito de términos.
\end{definicion}

\begin{ejemplo}[Fibonacci]
\[
f_n = f_{n-1} + f_{n-2}, \quad n \geq 3, \quad f_1 = f_2 = 1.
\]
\end{ejemplo}

\begin{ejemplo}[Torres de Hanoi]
\[
h_n = 2h_{n-1} + 1, \quad h_1 = 1.
\]
\end{ejemplo}

\section{Solución de Recurrencias}

\begin{definicion}
Resolver una relación de recurrencia significa hallar una fórmula explícita para $a_n$.
\end{definicion}

\begin{ejemplo}
$a_n = a_{n-1} + 3$, $a_1 = 2$ $\Rightarrow$ $a_n = 2 + 3(n-1)$.
\end{ejemplo}

\begin{ejemplo}
$h_n = 2h_{n-1} + 1$, $h_1 = 1$ $\Rightarrow$ $h_n = 2^n - 1$.
\end{ejemplo}

\begin{ejemplo}
$a_n = 3n\,a_{n-1}$, $a_0 = 2$ $\Rightarrow$ $a_n = 2 \cdot 3^n n!$.
\end{ejemplo}

\section{Recurrencias Homogéneas Lineales}

\begin{definicion}
Una recurrencia homogénea lineal de orden $k$ con coeficientes constantes es de la forma
\[
a_n = c_1 a_{n-1} + \cdots + c_k a_{n-k}, \quad c_k \neq 0.
\]
\end{definicion}

\begin{teorema}
Para la recurrencia de orden 2, $a_n = c_1 a_{n-1} + c_2 a_{n-2}$:
\begin{enumerate}
  \item Si $U_n, V_n$ son soluciones, $bU_n + dV_n$ también lo es.
  \item Si $r$ es raíz de $t^2 - c_1 t - c_2 = 0$, entonces $r^n$ es solución.
  \item Si $r_1 \neq r_2$ reales: $a_n = b r_1^n + d r_2^n$.
  \item Si $r$ es raíz doble: $a_n = b r^n + d n r^n$.
  \item Si las raíces son complejas: solución trigonométrica.
\end{enumerate}
\end{teorema}

\begin{ejemplo}[Fibonacci]
$f_n = f_{n-1} + f_{n-2}$, $f_0 = 0$, $f_1 = 1$. La ecuación característica es $r^2 - r - 1 = 0$, cuyas raíces son $\frac{1 \pm \sqrt{5}}{2}$. Por tanto,
\[
f_n = \frac{1}{\sqrt{5}}\left(\frac{1+\sqrt{5}}{2}\right)^n - \frac{1}{\sqrt{5}}\left(\frac{1-\sqrt{5}}{2}\right)^n.
\]
\end{ejemplo}

\begin{ejemplo}
$a_n = 6a_{n-1} - 9a_{n-2}$, $a_0 = 1$, $a_1 = -1$. La ecuación característica da $r = 3$ doble. Así $a_n = (b + dn)3^n$. Con las condiciones iniciales, $b = 1$, $d = -4/3$. Por tanto,
\[
a_n = \left(1 - \tfrac{4}{3}n\right)3^n.
\]
\end{ejemplo}

\section{Recurrencias No Homogéneas}

\begin{teorema}
Dada $a_n = c_1 a_{n-1} + c_2 a_{n-2} + f(n)$, si $g(n)$ es solución particular y $V_n$ es solución de la homogénea asociada, entonces $U_n = V_n + g(n)$ es solución general.
\end{teorema}

\begin{ejemplo}
$a_n = 6a_{n-1} - 8a_{n-2} + 3$, $a_0 = 0$, $a_1 = 1$.
Homogénea: $r^2 - 6r + 8 = 0 \Rightarrow r = 2, 4$.
Solución particular constante: $A = 6A - 8A + 3 \Rightarrow A = -1$.
Con las condiciones iniciales, $b = 0$, $d = 1/2$. Luego,
\[
a_n = \tfrac{1}{2}\,4^n - 1.
\]
\end{ejemplo}
\chapter{Funciones Generatrices}

\section{Funciones generatrices}

\begin{definicion}
La \textbf{función generatriz} de la sucesión $a_0, a_1, \ldots$ es
\[
f(x) = \sum_{i=0}^{\infty} a_i x^i.
\]
\end{definicion}

\begin{ejemplo}
$(1+x)^n$ es la función generatriz de $\binom{n}{0}, \binom{n}{1}, \ldots, \binom{n}{n}, 0, 0, \ldots$
\end{ejemplo}

\begin{ejemplo}
$\dfrac{1}{1-x} = 1 + x + x^2 + x^3 + \cdots$ es la función generatriz de $1, 1, 1, \ldots$
\end{ejemplo}

\begin{ejemplo}
$\dfrac{1}{(1-x)^2} = 1 + 2x + 3x^2 + \cdots$ es la función generatriz de $1, 2, 3, \ldots$
\end{ejemplo}

\begin{ejemplo}
La función generatriz de $0^2, 1^2, 2^2, \ldots$ es $\dfrac{x(1+x)}{(1-x)^3}$.
\end{ejemplo}

\begin{table}[h]
\centering
\caption{Funciones generatrices fundamentales}
\begin{tabular}{|c|l|}
\hline
\textbf{No.} & \textbf{Función generatriz} \\
\hline
1 & $(1+x)^m = \sum_{n=0}^m \binom{m}{n} x^n$ \\
2 & $(1+ax)^m = \sum_{n=0}^m \binom{m}{n} a^n x^n$ \\
3 & $(1+x^m)^n = \sum_{k=0}^n \binom{n}{k} x^{mk}$ \\
4 & $\frac{1-x^{m+1}}{1-x} = 1+x+\cdots+x^m$ \\
5 & $\frac{1}{1-x} = \sum_{n=0}^{\infty} x^n$ \\
6 & $\frac{1}{(1-x)^m} = \sum_{n=0}^{\infty} \binom{m+n-1}{n} x^n$ \\
7 & $\frac{1}{(1+x)^m} = \sum_{n=0}^{\infty} (-1)^n \binom{m+n-1}{n} x^n$ \\
8 & $e^x = \sum_{n=0}^{\infty} \frac{x^n}{n!}$ \\
9 & $\ln(1+x) = \sum_{n=1}^{\infty} (-1)^{n+1}\frac{x^n}{n}$ \\
10 & $\frac{1}{(1-x)^2} = \sum_{n=0}^{\infty} (n+1)x^n$ \\
\hline
\end{tabular}
\end{table}

\begin{ejemplo}
Mónica compró 12 naranjas para Graciela, María y Francisco. Graciela recibe al menos 4, María al menos 2 y Francisco entre 2 y 5. La función generatriz es:
\[
f(x) = (x^4 + \cdots + x^8)(x^2 + \cdots + x^6)(x^2 + \cdots + x^5).
\]
El coeficiente de $x^{12}$ es $14$.
\end{ejemplo}
\chapter{Análisis de Algoritmos}

% ============================================================
% Paquetes necesarios (ya deberían estar en el preámbulo)
% \usepackage{amsmath, amssymb, amsthm}
% \usepackage{listings, xcolor}
% \usepackage{tcolorbox}
% \tcbuselibrary{skins, breakable}
% \usepackage{tikz}
% \usepackage{pgfplots}
% \pgfplotsset{compat=1.18}
% ============================================================

% Colores para el código Python
\definecolor{pyblue}{RGB}{0,0,255}
\definecolor{pygreen}{RGB}{0,128,0}
\definecolor{pyred}{RGB}{255,0,0}
\definecolor{pygray}{RGB}{128,128,128}
\definecolor{pyback}{RGB}{245,245,245}

\lstset{
  language=Python,
  basicstyle=\small\ttfamily,
  keywordstyle=\color{pyblue}\bfseries,
  commentstyle=\color{pygreen}\itshape,
  stringstyle=\color{pyred},
  numberstyle=\tiny\color{pygray},
  stepnumber=1,
  numbers=left,
  backgroundcolor=\color{pyback},
  frame=single,
  breaklines=true,
  tabsize=4,
  showspaces=false,
  showstringspaces=false,
  captionpos=b,
}

% Estilos para cajas de teoremas
\newtcolorbox{definicionBox}{
  colback=blue!5,
  colframe=blue!75!black,
  fonttitle=\bfseries,
  title=Definición,
  breakable
}

\newtcolorbox{observacionBox}{
  colback=yellow!5,
  colframe=yellow!75!black,
  fonttitle=\bfseries,
  title=Observación,
  breakable
}

\newtcolorbox{ejemploBox}{
  colback=green!5,
  colframe=green!75!black,
  fonttitle=\bfseries,
  title=Ejemplo,
  breakable
}

% Redefinir entornos existentes (si ya están definidos, comentar)
% Si ya tienes definiciones con \newtheorem, puedes usar las cajas manualmente
% o redefinirlas. En este caso, mantendré los nombres originales pero con cajas.
% Para no romper la compilación, asumo que en el preámbulo ya están definidos
% los entornos definicion, observacion, ejemplo. Si no, los creo aquí.
% Para simplificar, usaré directamente los entornos con tcolorbox.


% Renombrar los entornos para usar cajas
\renewenvironment{definicion}
  {\begin{definicionBox}}
  {\end{definicionBox}}

\renewenvironment{observacion}
  {\begin{observacionBox}}
  {\end{observacionBox}}

\renewenvironment{ejemplo}
  {\begin{ejemploBox}}
  {\end{ejemploBox}}


% ============================================================
% INICIO DEL CAPÍTULO
% ============================================================

\section{Introducción}

\begin{definicion}
La \textbf{demora temporal} de un algoritmo (comúnmente conocida como tiempo de ejecución o complejidad temporal) es una medida que estima o calcula el tiempo que tarda un algoritmo en completarse en función del tamaño de sus datos de entrada.
\end{definicion}

\begin{observacion}
\begin{itemize}
  \item No se mide típicamente en segundos o minutos, sino en el número de operaciones fundamentales (o pasos) que realiza. Esto se debe a que el tiempo real en segundos depende de factores externos como la velocidad del hardware, el lenguaje de programación o la carga del sistema.
  \item El tiempo de ejecución depende de los datos de entrada. En lugar de analizar entradas específicas, se usan parámetros que definen su tamaño; por ejemplo, si la entrada es un conjunto que contiene $n$ elementos, se diría que el tamaño de la entrada es $n$.
\end{itemize}

\textbf{Tipos de análisis de tiempo:}
\begin{itemize}
  \item \textbf{Mejor caso:} tiempo mínimo necesario para una entrada de tamaño $n$.
  \item \textbf{Peor caso:} tiempo máximo necesario para una entrada de tamaño $n$.
  \item \textbf{Caso promedio:} tiempo necesario para un conjunto representativo de entradas de tamaño $n$.
\end{itemize}
\end{observacion}

\begin{observacion}[Estimación vs. cálculo exacto]
El objetivo es obtener una medida útil del tiempo, no un cálculo preciso. Para ello, se cuentan los pasos fundamentales o dominantes del algoritmo. Por ejemplo, si la actividad principal de un algoritmo es hacer comparaciones, como ocurre en una rutina para ordenar, se cuenta el número de comparaciones. Si un algoritmo consiste en un solo ciclo, se cuenta el número de iteraciones del ciclo.
\end{observacion}

\section{Notación Asintótica}

\begin{definicion}
Sean $T$ y $g$ dos funciones con dominio en los $\mathbb{Z}^+$. Se escribe:

\begin{itemize}
  \item $T(n) = O(g(n))$, y se dice que $T(n)$ es del orden a lo más $g(n)$, o $T(n)$ es $O$ mayúscula de $g(n)$, si existe una constante positiva $C_1$ tal que
  \[
  |T(n)| \leq C_1 |g(n)|
  \]
  para todos los enteros positivos $n$, excepto un número finito de ellos.

  \item $T(n) = \Omega(g(n))$, y se dice que $T(n)$ es del orden al menos $g(n)$, o $T(n)$ es omega de $g(n)$, si existe una constante positiva $C_2$ tal que
  \[
  |T(n)| \geq C_2 |g(n)|
  \]
  para todos los enteros positivos $n$, excepto un número finito de ellos.

  \item $T(n) = \Theta(g(n))$, y se dice que $T(n)$ es del orden $g(n)$, o $T(n)$ es theta de $g(n)$, si $T(n) = O(g(n))$ y $T(n) = \Omega(g(n))$.
\end{itemize}
\end{definicion}

\begin{observacion}
\begin{itemize}
  \item Si $T(n) = O(g(n))$ se dice que $g$ es una cota superior asintótica para $T$.
  \item Si $T(n) = \Omega(g(n))$ se dice que $g$ es una cota inferior asintótica para $T$.
  \item Si $T(n) = \Theta(g(n))$ se dice que $g$ es una cota estrecha asintótica para $T$.
\end{itemize}
\end{observacion}

\begin{ejemplo}[accesoPorIndice.py]
Encuentre una notación theta en términos de $n$, para el algoritmo de acceso por índice.

\begin{lstlisting}
def accesoPorIndice(lista, indice):
    return lista[indice]
\end{lstlisting}

$T(n) = 1$, pues cada elemento tiene una posición calculable, no se necesita recorrer la lista, solo se necesita una operación de memoria y es independiente del tamaño del dato de entrada. Por lo que, evidentemente:
\[
T(n) = \Theta(1)
\]
\end{ejemplo}

\begin{ejemplo}
Calcular el orden de la función $T(n) = 35n^2 + 7n + 4$.

\textbf{Solución.}
\begin{align*}
T(n) &= 35n^2 + 7n + 4 \leq 35n^2 + 7n^2 + 4n^2 = 46n^2 \\
T(n) &= O(n^2) \quad \text{(con $C_1 = 46$ en la Definición 4.4)}
\end{align*}
Además,
\[
T(n) = 35n^2 + 7n + 4 \geq 35n^2
\]
por lo que $T(n) = \Omega(n^2)$ (con $C_2 = 35$). De donde:
\[
T(n) = \Theta(n^2)
\]
\end{ejemplo}

\begin{observacion}
Es claro que el resultado del ejemplo anterior se puede generalizar: si
\[
T(n) = a_k n^k + a_{k-1} n^{k-1} + \cdots + a_1 n + a_0
\]
entonces:
\[
T(n) = \Theta(n^k)
\]
\end{observacion}

\begin{ejemplo}[busquedaLineal.py]
Encuentre una notación theta en términos de $n$, para el algoritmo de búsqueda lineal.

\begin{lstlisting}
def busquedaLineal(lista, elemento):
    n = len(lista)
    i = 0
    while i < n and lista[i] != elemento:
        i = i + 1
    if i < n:
        return i
    else:
        return print("elemento no encontrado")
\end{lstlisting}

\textbf{Mejor caso:} $T(n) = 1$, el elemento de la lista está al inicio.

\textbf{Peor caso:} $T(n) = n$, el elemento está al final de la lista o no se encuentra en la lista.

Entonces: $T(n) = \Theta(n)$.
\end{ejemplo}

\begin{observacion}
De aquí en adelante, los tiempos de demora los calcularemos en el peor de los casos.
\end{observacion}

\begin{ejemplo}[Multiplicación de matrices $n \times n$]
\begin{lstlisting}
def multiplicacion_matrices(A, B):
    n = len(A)
    # Crear matriz resultado llena de ceros
    C = [[0 for _ in range(n)] for _ in range(n)]
    for i in range(n):
        for j in range(n):
            for k in range(n):
                C[i][j] = A[i][k] * B[k][j]
    return C
\end{lstlisting}

\textbf{Solución.}
\[
T(n) = \underbrace{n^2 + n^2 + \cdots + n^2}_{n \text{ sumandos}} = n \cdot n^2 = n^3
\]
Entonces, $T(n) = \Theta(n^3)$.
\end{ejemplo}

\begin{ejemplo}
Encuentre una notación theta en términos de $n$, para el siguiente algoritmo, que cuenta el número total de iteraciones en un bucle anidado.

\begin{lstlisting}
def contarIteraciones(n):
    x = 0
    for i in range(1, n + 1):
        for j in range(1, i + 1):
            x += 1
    return x
\end{lstlisting}

\textbf{Solución.}
\[
T(n) = 1 + 2 + \cdots + n
\]
Para la cota superior:
\[
T(n) \leq n + n + \cdots + n = n^2 \quad \Rightarrow \quad T(n) = O(n^2)
\]
Para la cota inferior, no basta con $T(n) \ge n$; necesitamos una cota más ajustada:
\[
T(n) \ge \lceil n/2 \rceil + \cdots + n \ge \lceil n/2 \rceil \cdot \lceil n/2 \rceil \ge \frac{n^2}{4}
\]
por lo tanto $T(n) = \Omega(n^2)$. En consecuencia,
\[
T(n) = \Theta(n^2)
\]
\end{ejemplo}

\begin{ejemplo}
Encuentre una notación para el algoritmo que genera las combinaciones de tamaño $r$ de un conjunto de tamaño $n$.

\begin{lstlisting}
import math

def combinacion(n, r):
    # Inicializar la combinacion inicial [1, 2, 3, ..., r]
    s = [0] * (r + 1)  # Usamos indices 1..r
    for i in range(1, r + 1):
        s[i] = i

    # Imprimir la primera combinacion
    print([s[i] for i in range(1, r + 1)])

    # Calcular el numero total de combinaciones C(n, r)
    total_combinaciones = math.comb(n, r)

    # Generar las combinaciones restantes
    for count in range(2, total_combinaciones + 1):
        m = r
        max_val = n

        # Encontrar el elemento mas a la derecha que no tiene su valor maximo
        while s[m] == max_val:
            m -= 1
            max_val -= 1

        # Incrementar el elemento encontrado
        s[m] += 1

        # Los elementos a la derecha son sucesores del elemento incrementado
        for j in range(m + 1, r + 1):
            s[j] = s[j - 1] + 1

        # Imprimir la combinacion actual
        print([s[i] for i in range(1, r + 1)])
\end{lstlisting}

\textbf{Solución.}
\[
T(n) = 1^r + 2^r + \cdots + n^r \leq n \cdot n^r = n^{r+1}
\]
luego $T(n) = O(n^{r+1})$.

Por otro lado,
\[
T(n) \ge \lceil n/2 \rceil^r + \cdots + n^r \ge \frac{n}{2} \cdot \left(\frac{n}{2}\right)^r = \frac{n^{r+1}}{2^{r+1}}
\]
por lo tanto $T(n) = \Omega(n^{r+1})$. En consecuencia,
\[
T(n) = \Theta(n^{r+1})
\]
\end{ejemplo}

\begin{ejemplo}
Encuentre una notación theta en términos de $n$, para el siguiente algoritmo.

\begin{lstlisting}
def generaSubsets(elements):
    # elements = lista de elementos
    n = len(elements)
    subsets = []
    for i in range(1 << n):  # 2^n iteraciones
        subset = [elements[j] for j in range(n) if (i >> j) & 1]
        subsets.append(subset)
    return subsets
\end{lstlisting}

\textbf{Solución.}
Cada iteración del bucle genera un subconjunto, y la comprensión interna recorre $n$ elementos, por lo que el costo total es
\[
T(n) = n \cdot 2^n
\]
Luego,
\[
T(n) = \Theta(n \cdot 2^n)
\]
\end{ejemplo}

\begin{observacion}
La clave en la práctica es siempre hacer esfuerzos para encontrar soluciones con las complejidades más bajas posibles (idealmente, $\Theta(1)$, $\Theta(\log n)$, $\Theta(n)$ o $\Theta(n \log n)$). Entender estas funciones es crucial para:

\begin{itemize}
  \item \textbf{Elegir el algoritmo correcto:} saber si un problema necesita una solución $\Theta(n \log n)$ o si te puedes permitir una $\Theta(n^2)$.
  \item \textbf{Predecir el rendimiento:} anticipar cómo se comportará un algoritmo si los datos de entrada crecen.
  \item \textbf{Identificar cuellos de botella:} saber que un algoritmo $\Theta(n^3)$ dentro de tu código es probablemente el causante de la lentitud.
  \item Saber que algoritmos con $\Theta(2^n)$, $\Theta(n!)$, $\Theta(k^n)$ con $k > 1$ son intratables para entradas de tamaño moderado. Solo son útiles para datos de entrada muy pequeños.
\end{itemize}

En la Figura~\ref{fig:crecimiento} puede visualizarse cómo es el crecimiento de estas funciones.
\end{observacion}

\begin{figure}[htbp]
\centering
\begin{tikzpicture}
\begin{axis}[
    width=0.85\textwidth,
    height=7cm,
    xlabel={$n$},
    ylabel={$y$},
    xmin=1, xmax=13,
    ymin=1, ymax=300,
    ymode=log,
    log basis y=2,
    ytick={1,2,4,8,16,32,64,128,256},
    xtick={1,2,3,4,5,6,7,8,9,10,11,12,13},
    axis lines=left,
    legend pos=outer north east,
    legend style={font=\small},
    every axis plot/.append style={thick},
]
\addplot[domain=1:9, samples=100, black, thick] {2^x};
\addlegendentry{$y = 2^n$}

\addplot[domain=1:13, samples=100, black!70, thick] {x^2};
\addlegendentry{$y = n^2$}

\addplot[domain=1:13, samples=100, gray, thick] {x*ln(x)/ln(2)};
\addlegendentry{$y = n \lg n$}

\addplot[domain=1:13, samples=100, black, thin] {x};
\addlegendentry{$y = n$}

\addplot[domain=1:13, samples=100, gray, thin] {ln(x)/ln(2)};
\addlegendentry{$y = \lg n$}

\addplot[domain=1:13, samples=2, black, dashed] {1};
\addlegendentry{$y = 1$}
\end{axis}
\end{tikzpicture}
\caption{Crecimiento de algunas funciones comunes en la complejidad computacional.}
\label{fig:crecimiento}
\end{figure}

\begin{ejemplo}
Encuentre una notación $\Theta$ para el número de veces que se ejecuta la asignación $x = x + 1$ en el siguiente algoritmo.

\begin{lstlisting}
j = n
while (j >= 1) {
    for i = 1 to j
        x = x + 1
    j = floor(j/2)
}
\end{lstlisting}

\textbf{Solución.}
\[
T(n) = n + \left\lfloor \frac{n}{2} \right\rfloor + \left\lfloor \frac{n}{4} \right\rfloor + \cdots + \left\lfloor \frac{n}{2^{k-1}} \right\rfloor
\]
donde $k$ denota el número de veces que se ejecuta el cuerpo del ciclo \texttt{while}, es decir, el número de veces que se ejecuta $x = x + 1$.

Acotando superiormente:
\[
T(n) \leq n + \frac{n}{2} + \frac{n}{4} + \cdots + \frac{n}{2^{k-1}}
= \frac{n\left(1 - \frac{1}{2^k}\right)}{1 - \frac{1}{2}}
= 2n\left(1 - \frac{1}{2^k}\right) \leq 2n
\]
luego $T(n) = O(n)$. Además, como cada término es positivo,
\[
T(n) \ge n \quad \Rightarrow \quad T(n) = \Omega(n)
\]
Por lo tanto,
\[
T(n) = \Theta(n)
\]
\end{ejemplo}

\begin{observacion}
Un problema se considera \textbf{tratable} o factible cuando existe un algoritmo que lo resuelve con un tiempo de ejecución polinomial en el peor de los casos. Esta condición se interpreta como la posesión de una solución eficiente. No obstante, si el tiempo requerido es proporcional a un polinomio de grado elevado, la solución podría ser ineficiente en la práctica. Afortunadamente, en muchos casos de importancia, el grado del polinomio que acota el tiempo de ejecución es bajo, lo que garantiza una verdadera eficiencia.

Se denomina \textbf{intratable} a un problema para el cual no se conoce ningún algoritmo que lo resuelva en tiempo polinomial en el peor caso. Esto implica que cualquier algoritmo que lo solucione requerirá un tiempo de ejecución prohibitivamente largo, generalmente de naturaleza exponencial o factorial, incluso para instancias o entradas de tamaño moderado.

Se dice que un problema es \textbf{irresoluble} o no computable cuando no existe ningún algoritmo capaz de resolverlo. Un ejemplo fundamental, y uno de los primeros en ser demostrado como tal, es el problema de la parada\footnote{Halting problem.}: determinar si un programa cualquiera se detendrá en algún momento, dada una entrada específica. Existe una variedad significativa de problemas de este tipo, varios de los cuales revisten gran importancia práctica en el campo de la computación teórica.

Existe una gran cantidad de problemas computacionales solucionables cuyo estatus de complejidad permanece indeterminado; si bien se conjetura que son intratables, aún no se ha demostrado formalmente que lo sean. La mayoría de estos problemas pertenecen a la clase de los NP-completos.
\end{observacion}

\part{Teoría de Grafos}
\chapter{Teoría de Grafos}

\section{Conceptos fundamentales}

\subsection{Definiciones básicas}

\begin{definicion}
Una \textbf{familia enumerable} es una colección enumerable de objetos iguales o diferentes.
\begin{itemize}
  \item El mayor número $p$ de elementos iguales en una familia enumerable le llamamos \textbf{grado de multiplicidad}.
  \item Una familia enumerable con $p = 1$ se llama \textbf{conjunto}.
  \item En lo que sigue, consideraremos solo familias finitas.
  \item El conjunto de partes de $X$ será denotado $\mathcal{P}(X)$.
  \item $\mathcal{P}_k(X) = \{U \in \mathcal{P}(X) \mid |U| = k\}$.
  \item La potencia cartesiana $X^k$ es el conjunto de todas las $k$-uplas ordenadas $(x_1, x_2, \ldots, x_k)$ donde $x_i \in X$, $1 \leq i \leq k$, $1 \leq k \leq |X|$.
\end{itemize}
\end{definicion}

\begin{definicion}
Un \textbf{grafo} es una estructura $G = (N, A)$ donde $N$ es un conjunto finito y $A$ es una familia cuyos elementos (no vacíos) son definidos en función de los elementos $v$ de $N$, en dos formas posibles:

\begin{enumerate}
  \item Si la familia $A$ de grado de multiplicidad $p$ fuera tal que sus elementos $a \in \mathcal{P}(N)$, $G$ será llamado un \textbf{$p$-hipergrafo no orientado}.
  \item Si la familia $A$ de grado de multiplicidad $p$ es constituida por pares de elementos no vacíos de $\mathcal{P}(N)$ de la forma $a_i = (a_i^-, a_i^+)$, $G$ será llamado un \textbf{$p$-hipergrafo orientado}.
\end{enumerate}

\begin{itemize}
  \item Los elementos de $N$ son llamados \textbf{nodos} (o vértices) y el valor $n = |N|$ es el \textbf{orden} de la estructura.
  \item La familia $A$ puede ser entendida como una relación o conjunto de relaciones de adyacencia cuyos elementos son llamados en general \textbf{conexiones}; en particular, en las estructuras no orientadas las $a \in A$ son conocidas como \textbf{aristas} y en las orientadas como \textbf{arcos}.
  \item Dos nodos que participan de una conexión son llamados \textbf{adyacentes}.
  \item Al valor $m = |A|$ lo llamaremos el \textbf{tamaño} del grafo.
  \item Si $m = 0$, el grafo es llamado \textbf{trivial}.
  \item La representación dinámica de un grafo, también llamado grafo (confusión semiintensional), es obtenida asociando a cada nodo un punto o una pequeña área delimitada por una frontera, y a cada conexión un diseño capaz de representar la forma de asociación de los nodos que involucra la conexión.
\end{itemize}
\end{definicion}

\begin{ejemplo}
Los siguientes gráficos representan aristas o arcos. (Ver Figura~\ref{fig:aristas-arcos}.)
\end{ejemplo}

\begin{figure}[htbp]
\centering
\begin{tikzpicture}[scale=0.9, every node/.style={circle, fill=black, inner sep=1.5pt}]
% Arista
\begin{scope}[xshift=0cm]
\node[label=left:{1}] (a1) at (0,0) {};
\node[label=right:{2}] (a2) at (0.8,1.2) {};
\draw[thick] (a1) to[bend left=20] (a2);
\node[fill=none] at (0.4,-0.7) {Arista $\{1,2\}$};
\end{scope}

% Arco
\begin{scope}[xshift=3cm]
\node[label=left:{1}] (b1) at (0,0) {};
\node[label=right:{2}] (b2) at (0.8,1.2) {};
\draw[thick,->] (b1) to[bend left=20] (b2);
\node[fill=none] at (0.4,-0.7) {Arco $(1,2)$};
\end{scope}

% Hiper-arista
\begin{scope}[xshift=6.5cm]
\node[label=above left:{1}] (c1) at (0,1) {};
\node[label=above right:{2}] (c2) at (1.2,1) {};
\node[label=below left:{4}] (c4) at (0,0) {};
\node[label=below right:{3}] (c3) at (1.2,0) {};
\draw[thick] plot [smooth cycle, tension=0.7] coordinates {(-0.35,1.25) (1.3,1.3) (1.35,-0.3) (-0.3,-0.3)};
\node[fill=none] at (0.6,-0.9) {Hiper-arista $\{1,2,3\}$};
\end{scope}

% Hiper-arco
\begin{scope}[xshift=10.5cm]
\node[label=above left:{1}] (d1) at (0,1) {};
\node[label=above right:{2}] (d2) at (1.2,1) {};
\node[label=below left:{4}] (d4) at (0,0) {};
\node[label=below right:{3}] (d3) at (1.2,0) {};
\draw[thick,->] (d1) to[bend right=15] (d3);
\draw[thick,->] (d4) to[bend left=15] (d2);
\node[fill=none] at (0.6,-0.9) {Hiper-arco $(\{4,1\},\{2,3\})$};
\end{scope}
\end{tikzpicture}
\caption{Utilizaremos el término grafo para designar estructuras que posean conexiones con un máximo de dos nodos involucrados y, a menos que se indique lo contrario, con grado de multiplicidad igual a 1.}
\label{fig:aristas-arcos}
\end{figure}

\begin{observacion}
\begin{enumerate}
  \item[a)] Para $G = (N, A)$ un $p$-hipergrafo no orientado, en particular, si:
  \begin{itemize}
    \item Se verifica que todos los elementos de $A$ pertenecen a $\mathcal{P}_1(N) \cup \mathcal{P}_2(N)$; $G$ es llamado un \textbf{multigrafo} o \textbf{$p$-grafo no orientado}.
    \item $A$ fuese un subconjunto de $\mathcal{P}(N)$, $G$ será \textbf{1-hipergrafo no orientado} o \textbf{hipergrafo simple}.
  \end{itemize}
  \item[b)] Para $G = (N, A)$ un 1-hipergrafo orientado, en particular, si:
  \begin{itemize}
    \item Se tuviera $|a_i^-| = |a_i^+| = 1$ para todo $i$, $G$ se denominará \textbf{1-grafo orientado} o \textbf{dígrafo}.
  \end{itemize}
  \item[c)] Una conexión que tiene apenas un nodo es llamada \textbf{bucle} o \textbf{lazo}. En las estructuras no orientadas un lazo es un elemento de $\mathcal{P}_1(N)$ y en las orientadas, una $k$-upla de elementos iguales; en grafos orientados, un par ordenado de elementos iguales.
  \item[d)] Asociado a todo grafo orientado $G = (N, A)$ existe un grafo $G' = (N, A')$ en el cual todo arco $a' = (y, x)$ tiene en $A$ un arco correspondiente $a = (x, y)$. Al grafo $G'$ se le denomina el \textbf{dual direccional} de $G$.
\end{enumerate}
\end{observacion}

\begin{ejemplo}
A continuación se presentan diversos tipos de estructuras de grafo: 3-hipergrafo no orientado, 1-hipergrafo orientado, multigrafo, grafo orientado, hipergrafo simple, grafo simple, 2-grafo orientado y grafo orientado con un lazo.
\end{ejemplo}

\begin{definicion}
Dos grafos $G_1 = (N_1, A_1)$ y $G_2 = (N_2, A_2)$ son:
\begin{enumerate}
  \item[a)] \textbf{Iguales} cuando $N_1 = N_2$ y $A_1 = A_2$.
  \item[b)] \textbf{Isomorfos} cuando existe una biyección $f : N_1 \to N_2$ que preserve las relaciones de adyacencia, es decir, que para todos $x, y \in N_1$, $(x,y) \in A_1$ ($\{x,y\} \in A_1$) si y solo si $(f(x),f(y)) \in A_2$ ($\{f(x),f(y)\} \in A_2$). A la función $f$ se le denomina \textbf{isomorfismo}.
\end{enumerate}
\end{definicion}

\begin{ejemplo}
Los grafos $G_1$ y $G_2$ son isomorfos, lo que podemos denotar con $G_1 \cong G_2$. Ver Figura~\ref{fig:isomorfismo}.
\end{ejemplo}

\begin{figure}[htbp]
\centering
\begin{tikzpicture}[scale=1, every node/.style={circle, fill=black, inner sep=1.5pt}]
% G1: pentagon
\begin{scope}
\node[label=above:{a}] (a) at (0,2) {};
\node[label=right:{b}] (b) at (1.3,0.8) {};
\node[label=below right:{c}] (c) at (0.8,-0.8) {};
\node[label=below left:{d}] (d) at (-0.8,-0.8) {};
\node[label=left:{e}] (e) at (-1.3,0.8) {};
\draw (a)--(b)--(c)--(d)--(e)--(a);
\node[fill=none] at (0,-1.6) {$G_1$};
\end{scope}

% G2: pentagram
\begin{scope}[xshift=6cm]
\node[label=above:{1}] (n1) at (0,2) {};
\node[label=right:{4}] (n4) at (1.3,0.8) {};
\node[label=below right:{2}] (n2) at (0.8,-0.8) {};
\node[label=below left:{5}] (n5) at (-0.8,-0.8) {};
\node[label=left:{3}] (n3) at (-1.3,0.8) {};
\draw (n1)--(n2)--(n3)--(n4)--(n5)--(n1);
\node[fill=none] at (0,-1.6) {$G_2$};
\end{scope}
\end{tikzpicture}
\caption{La función $f: N_1 \longrightarrow N_2$ definida por $f(a)=1$, $f(b)=2$, $f(c)=3$, $f(d)=4$, $f(e)=5$ es un isomorfismo.}
\label{fig:isomorfismo}
\end{figure}

\begin{observacion}
\begin{enumerate}
  \item[a)] Diremos que un grafo es \textbf{valorado} sobre los nodos (conexiones) cuando existen una o más funciones relacionando $N$ ($A$) con un conjunto de números.
  \item[b)] En la mayoría de las aplicaciones de grafos existen datos cuantitativos asociados a nodos o conexiones involucradas en el problema; los modelos correspondientes involucran, en esos casos, grafos valorados.
\end{enumerate}
\end{observacion}

\begin{ejemplo}
Los turistas Jenssen, Leuzinger, Dufour y Medeiros se encuentran en un bar de París y comienzan a conversar. Las lenguas disponibles son el inglés, el francés, el portugués y el alemán; Jenssen habla todas, Leuzinger no habla apenas el portugués, Dufour habla francés y alemán, y Medeiros habla inglés y portugués.

\begin{enumerate}
  \item[a)] Represente por medio de un grafo $G = (N, A)$ todas las posibilidades de que uno de ellos dirija la palabra a otro, siendo comprendido. Defina $N$ y $A$.
  \item[b)] Represente por medio de un hipergrafo $H = (N, A)$ las capacidades lingüísticas del grupo. ¿Cuál es el significado de las intersecciones $a_i \cap a_j$ donde los $a_k \in A$?
\end{enumerate}
\end{ejemplo}

\begin{definicion}
Un grafo $G = (N, A)$ se denomina \textbf{grafo $k$-partido} si $N$ se puede particionar en subconjuntos $N_1, N_2, \ldots, N_k$, tales que no haya dos nodos adyacentes en un mismo $N_i$.
\end{definicion}

\begin{ejemplo}
La Figura~\ref{fig:bipartido} muestra un grafo bipartido y un grafo no bipartido.
\end{ejemplo}

\begin{figure}[htbp]
\centering
\begin{tikzpicture}[scale=1, every node/.style={circle, fill=black, inner sep=1.5pt}]
\begin{scope}
\node[label=left:{a}] (a) at (0,1.5) {};
\node[label=left:{b}] (b) at (0,0) {};
\node[label=right:{c}] (c) at (2,1.5) {};
\node[label=right:{d}] (d) at (2,0.5) {};
\node[label=right:{e}] (e) at (2,-0.5) {};
\draw (a)--(c);
\draw (a)--(d);
\draw (b)--(c);
\draw (b)--(d);
\draw (b)--(e);
\node[fill=none] at (1,-1.2) {Bipartido};
\end{scope}

\begin{scope}[xshift=5cm]
\node[label=above:{a}] (a2) at (0.7,1.5) {};
\node[label=left:{c}] (c2) at (0,0) {};
\node[label=right:{b}] (b2) at (1.4,0) {};
\draw (a2)--(c2)--(b2)--(a2);
\node[fill=none] at (0.7,-0.9) {No bipartido};
\end{scope}
\end{tikzpicture}
\caption{Cuando $k=2$ el grafo $G$ es llamado \textbf{bipartido}.}
\label{fig:bipartido}
\end{figure}

\begin{ejemplo}
La siguiente figura muestra un grafo 3-partido (nodos $a,b$; $c,d$; $e,f$).
\end{ejemplo}

\begin{ejemplo}[Los puentes de Königsberg]
Durante el siglo \textsc{xvii}, la ciudad de Königsberg (antes en Prusia Oriental; llamada ahora Kaliningrado, en Rusia) estaba dividida en cuatro zonas (incluida la isla de Kneiphof) por el río Pregel. Siete puentes comunicaban estas regiones. Se decía que los habitantes hacían paseos dominicales tratando de encontrar una forma de caminar por la ciudad cruzando cada puente exactamente una vez y regresando al punto donde se había iniciado el paseo. Con el fin de determinar si existía o no dicho circuito, en 1736 el matemático suizo Leonhard Euler representó las cuatro zonas de la ciudad y los siete puentes con el multigrafo correspondiente (ver Figura~\ref{fig:konigsberg}), dando así inicio a la Teoría de Grafos.
\end{ejemplo}

\begin{figure}[htbp]
\centering
\begin{tikzpicture}[scale=1, every node/.style={circle, fill=black, inner sep=1.5pt}]
\node[label=above:{c}] (c) at (0,1.5) {};
\node[label=left:{a}] (a) at (0,0) {};
\node[label=below:{b}] (b) at (0,-1.5) {};
\node[label=right:{d}] (d) at (2.5,0) {};
\draw (a) to[bend left=25] (c);
\draw (a) to[bend right=25] (c);
\draw (a) to[bend left=25] (b);
\draw (a) to[bend right=25] (b);
\draw (c) -- (d);
\draw (a) -- (d);
\draw (b) -- (d);
\end{tikzpicture}
\caption{Multigrafo que representa los 7 puentes de Königsberg: cuatro zonas (nodos $a,b,c,d$) y siete puentes (aristas).}
\label{fig:konigsberg}
\end{figure}

\section{Representación numérica de los grafos}

Aunque la representación visual de un grafo resulta accesible para el análisis humano, no es viable para su implementación en sistemas computacionales. En consecuencia, se hace imprescindible recurrir a representaciones numéricas que posibiliten el tratamiento automatizado de la información estructural. La literatura ofrece múltiples alternativas para la representación numérica de grafos, cuya idoneidad varía en función del contexto de aplicación. Dichas representaciones pueden estar orientadas a optimizar operaciones de búsqueda y almacenamiento (enfoque algorítmico) o a facilitar el desarrollo de razonamientos de carácter algebraico o combinatorio.

\subsection{Matriz de adyacencia}

Las diversas representaciones matriciales de estructuras de grafo están habitualmente asociadas a la necesidad de realización de cálculos que involucran datos estructurales.

Dado un grafo orientado $G$ con $n$ nodos numerados $v_1, v_2, \ldots, v_n$ (esta numeración define una ordenación arbitraria en el conjunto de nodos) se puede definir la matriz $A$ de orden $n \times n$, donde:
\[
[A]_{ij} =
\begin{cases}
1 & \text{si } (v_i, v_j) \in A \\
0 & \text{si } (v_i, v_j) \notin A
\end{cases}
\]
que denominaremos \textbf{matriz de adyacencia}.

\begin{observacion}
\begin{itemize}
  \item Si un grafo orientado $G$ tiene matriz de adyacencia $A$ entonces $A^T$ es la matriz de adyacencia del dual direccional de $G$.
  \item La matriz de adyacencia de un grafo no orientado se puede representar en su forma triangular (superior o inferior), en vista que no hay distinción entre los dos sentidos posibles para cada arista.
\end{itemize}
\end{observacion}

\begin{ejemplo}
Un dígrafo de 5 nodos con lazos en 1 y 2 tiene matriz de adyacencia:
\[
A =
\begin{bmatrix}
1 & 0 & 0 & 0 & 1 \\
0 & 1 & 1 & 0 & 0 \\
0 & 0 & 0 & 0 & 1 \\
1 & 1 & 0 & 0 & 0 \\
0 & 1 & 0 & 0 & 0
\end{bmatrix}
\]

Un grafo simple de 4 nodos, representado en forma triangular:
\[
A =
\begin{bmatrix}
1 & 0 & 0 & 0 \\
1 & 0 & 0 & 0 \\
0 & 1 & 0 & 0 \\
1 & 0 & 1 & 0
\end{bmatrix}
\]

Un 3-hipergrafo no orientado de 4 nodos, con multiplicidad:
\[
A =
\begin{bmatrix}
1 & 0 & 0 & 0 \\
1 & 0 & 0 & 0 \\
0 & 3 & 0 & 0 \\
0 & 1 & 1 & 0
\end{bmatrix}
\]
\end{ejemplo}

Como limitaciones, cabe observar que esta representación no es adecuada para hipergrafos en general. Con $p$-grafos podrá también haber dificultades: en este caso los valores atribuidos a las posiciones no nulas serán los órdenes de multiplicidad de los arcos o aristas, y la matriz, por tanto, solamente sería formada apenas por valores $0$ y $1$. Tenga presente que mucha de la teoría está basada en la posibilidad de ver la matriz de adyacencia como matriz booleana.

Otra de las desventajas de usar matriz de adyacencia para representar un grafo se debe a que este tipo de representación hace difícil almacenar información adicional acerca del grafo; por ejemplo, si es preciso incluir información acerca de los nodos, será preciso representarla mediante alguna estructura adicional de almacenamiento.

\subsection{Lista de adyacencia}

Esta forma es la más conveniente para la entrada de datos por la simplicidad y ``economía'' de su representación. Es construida como un conjunto de listas de nodos; cada lista está formada por un nodo y por un conjunto de nodos que reciben de él un arco o que con él comparten una arista.

La matriz de adyacencia no es la más ``económica'' de las formas de representación de grafos; ella ocupa $n^2$ posiciones en memoria, en cuanto la lista de adyacencia utiliza apenas $n + m$ posiciones.

\begin{ejemplo}
Para un dígrafo con nodos $\{a,b,c,d,e,f\}$:

\begin{center}
\begin{tabular}{cc}
\begin{tabular}{c|c}
\multicolumn{2}{c}{Nodos} \\
Origen & Destino \\ \hline
a & b \\
b & a c \\
c & a d e f \\
d & a f \\
e & a \\
f & c e
\end{tabular}
&
\begin{tabular}{c|c}
\multicolumn{2}{c}{Nodos} \\
Destino & Origen \\ \hline
a & b c d \\
b & a \\
c & b f \\
d & c \\
e & c f \\
f & c d
\end{tabular}
\end{tabular}
\end{center}
\end{ejemplo}

\section{Relaciones de Adyacencia}

\begin{definicion}
En un grafo orientado $G = (N, A)$ se dice que $y \in N$ es \textbf{sucesor} de $x \in N$, cuando existe $(x,y) \in A$. Una definición dual direccional correspondiente es la de \textbf{antecesor}; luego, habiendo $(x,y) \in A$, $x$ es antecesor de $y$.
\end{definicion}

\begin{definicion}
\textbf{Vecino} o nodo adyacente de un nodo $x$, en un grafo orientado o no, es todo nodo que participa de una conexión con $x$.
\end{definicion}

Los conjuntos de sucesores y de antecesores de un nodo $x$ son denotados, respectivamente, $\Gamma^+(x)$ y $\Gamma^-(x)$, y corresponden a los conjuntos de elementos no nulos de la fila y de la columna asociadas a $x$ en la matriz de adyacencia del grafo. El conjunto de los vecinos de $x \in N$ es denotado $\Gamma(x)$.

\begin{observacion}
\begin{itemize}
  \item La información contenida en los conjuntos de sucesores o antecesores, o de vecinos, equivale a la contenida en el conjunto de conexiones. Por tanto, se puede representar un 1-grafo orientado como $G = (N, \Gamma^+)$ o $G = (N, \Gamma^-)$, y un 1-grafo no orientado por $G = (N, \Gamma)$.
  \item Estos conceptos pueden ser extendidos a subconjuntos de nodos; por ejemplo, si tenemos que $A \subseteq N$:
  \[
  \Gamma^+(A) = \bigcup_{x \in A} \Gamma^+(x)
  \]
  \item Cabe observar que la notación utilizada es abreviada; tratándose, de hecho, de componentes entre conjuntos, la notación correcta es del tipo $\Gamma(\{x\})$.
  \item Las nociones de sucesor y de antecesor pueden ser aplicadas iterativamente, llevando a conceptos que implican la conexión directa o indirecta entre nodos en grafos orientados. Los conjuntos así obtenidos son conocidos como \textbf{clausuras transitivas}.
\end{itemize}
\end{observacion}

\begin{definicion}
Un nodo $y$ es \textbf{alcanzable} a partir de un nodo $x$, cuando existe en el grafo considerado una secuencia de sucesores que se inicia en $x$ y llega a $y$.
\end{definicion}

\begin{definicion}
La \textbf{clausura transitiva directa} $\hat{\Gamma}^+(x)$ de un nodo $x$ es el conjunto de los nodos alcanzables a partir de $x$, y la \textbf{clausura transitiva inversa} $\hat{\Gamma}^-(x)$ de $x$ es el conjunto de los nodos o partes de los cuales $x$ es alcanzable.
\end{definicion}

Estas dos definiciones de clausura son duales direccionales. Es claro que, si $y \in \hat{\Gamma}^+(x)$, $y$ es un \textbf{descendiente} de $x$; si $y \in \hat{\Gamma}^-(x)$, $y$ es un \textbf{ascendiente} de $x$.

\begin{observacion}
Presentamos a continuación el proceso iterativo que conduce a la determinación de $\hat{\Gamma}^+(x)$. Por coherencia con la idea de alcanzable, se acepta al nodo origen como alcanzable a partir de cero etapas:
\[
\Gamma^0(x) = \{x\}, \quad \Gamma^{+1}(x) = \Gamma^+(x), \quad \Gamma^{+2}(x) = \Gamma^+(\Gamma^{+1}(x)), \quad \ldots, \quad \Gamma^{+n}(x) = \Gamma^+(\Gamma^{+n-1}(x))
\]
\[
\hat{\Gamma}^+(x) = \bigcup_{k=0}^{n} \Gamma^{+k}(x)
\]

La noción de clausura transitiva está conceptualmente asociada a las ideas de comunicación y de control. Si definimos un grafo $G = (N, A)$ considerando $N = $ Personas y $A = \{(x,y) \mid x \text{ conoce a } y\}$, la clausura transitiva directa de una persona $k$ será el conjunto de todas las personas que podrán tomar conocimiento de una información que $k$ disponga, a través de comunicación de persona a persona (o sea, sin el uso de los medios de comunicación de masa). La noción puede ser usada, por tanto, para estudiar el efecto de la comunicación informal conocida como ``tele-tipo''.
\end{observacion}

\section{Vecindades de conexiones}

Las definiciones que a continuación se dan son, explícitamente, para grafos sin lazos.

\begin{definicion}
\begin{enumerate}
  \item[a)] Dos conexiones son \textbf{adyacentes} si tienen un nodo en común.
  \item[b)] Se dice que una conexión es \textbf{incidente} en un nodo cuando el nodo forma parte de la conexión. En el caso de grafos orientados, un arco incide exteriormente en $x \in N$ si $x$ es el nodo del cual parte, e incide interiormente en $x$ si $x$ es el nodo al cual llega.
  \item[c)] Para un conjunto de nodos $A \subseteq N$, decimos que un arco incide exterior o interiormente en $A$ si incide de la misma manera en al menos un nodo de $A$.
\end{enumerate}
\end{definicion}

\begin{observacion}
En el caso de grafos no orientados, simplemente se dice que una arista incide en $A \subseteq N$. El conjunto de arcos incidentes exteriormente en $A$ se denota por $\omega^+(A)$, mientras que el conjunto de arcos incidentes interiormente se denota por $\omega^-(A)$; sus cardinalidades son, respectivamente, el \textbf{semigrado exterior} $\delta^+(A)$ y el \textbf{semigrado interior} $\delta^-(A)$.

Para grafos no orientados, el conjunto de aristas incidentes en $A$ se denota simplemente por $\omega(A)$.
\end{observacion}

\begin{ejemplo}
Considere el siguiente grafo no orientado $G = (N,A)$ con nodos $\{1,2,3,4\}$ y aristas $\{1,2\}, \{1,3\}, \{2,4\}, \{3,4\}$.

Tomemos el conjunto de nodos $A = \{1,2\} \subset N$. Las aristas incidentes en $A$ son aquellas que tienen al menos un extremo en $A$:
\[
\omega(A) = \{\{1,2\}, \{1,3\}, \{2,4\}\}
\]
Observe que la arista $\{3,4\}$ no está en $\omega(A)$ porque ninguno de sus extremos pertenece a $A$.

Si el grafo fuera orientado, con arcos $(1,2), (1,3), (2,4), (4,2)$: para el mismo conjunto $A = \{1,2\}$, los conjuntos de arcos incidentes son:
\[
\omega^+(A) = \{(1,2), (1,3)\}
\]
pues estos arcos tienen su extremo inicial en $A$, y
\[
\omega^-(A) = \{(1,2), (4,2)\}
\]
pues estos arcos tienen su extremo final en $A$.

Sus cardinalidades son:
\[
\delta^+(A) = |\omega^+(A)| = 2, \qquad \delta^-(A) = |\omega^-(A)| = 2
\]
Es decir, el semigrado exterior de $A$ es 2 y el semigrado interior de $A$ también es 2.
\end{ejemplo}

\begin{definicion}
El \textbf{grado} $\delta(x)$ de un nodo es el número total de conexiones que en él inciden.
\end{definicion}

\begin{observacion}
\begin{itemize}
  \item En grafos orientados, tenemos evidentemente que $\delta(x) = \delta^+(x) + \delta^-(x)$, y en grafos no orientados $\delta(x) = |\omega(x)|$. Se trata de una definición no orientada válida para $p$-grafos en general.
  \item Un grafo no orientado en el cual $\delta(x) = k$ para todo $x \in N$ es llamado \textbf{regular} (de grado $k$).
  \item Llamamos a un grafo orientado en el cual $\delta^+(x) = k$ ($\delta^-(x) = k$) para todo $x \in N$, \textbf{exteriormente regular} (\textbf{interiormente regular}) (de semigrado $k$).
  \item Para 1-grafos, las cardinalidades de los conjuntos de sucesores (antecesores) de un nodo son iguales a los conjuntos de conexiones exteriormente (interiormente) incidentes que les corresponden.
\end{itemize}
\end{observacion}

\begin{definicion}
Un grafo $G = (N, A)$ será \textbf{completo} si existe al menos una conexión asociada a cada par de nodos.
\end{definicion}

\begin{observacion}
\begin{itemize}
  \item En el caso no orientado significa exactamente una conexión.
  \item El grafo completo poseerá todos los nodos, es decir $G = (N, \mathcal{P}_2(N))$. Evidentemente, todas las estructuras de ese tipo con un mismo orden serán isomorfas; ellas son conocidas como \textbf{grafos completos} y se denotan $K_n$.
  \item Los grafos bipartidos no orientados con el mayor número posible de aristas son llamados \textbf{grafos bipartidos completos} y se denotan con $K_{p,q}$, y el número de aristas será, por tanto, $p \cdot q$.
\end{itemize}
\end{observacion}

\begin{ejemplo}
Los grafos completos $K_1$ a $K_5$ se muestran en la Figura~\ref{fig:completos}.
\end{ejemplo}

\begin{figure}[htbp]
\centering
\begin{tikzpicture}[scale=0.8, every node/.style={circle, fill=black, inner sep=1.5pt}]
\begin{scope}
\node (k1) at (0,0) {};
\node[fill=none] at (0,-0.7) {$K_1$};
\end{scope}

\begin{scope}[xshift=2cm]
\node (a) at (0,0) {};
\node (b) at (1,0) {};
\draw (a)--(b);
\node[fill=none] at (0.5,-0.7) {$K_2$};
\end{scope}

\begin{scope}[xshift=5cm]
\node (a) at (0,0) {};
\node (b) at (1,0) {};
\node (c) at (0.5,0.9) {};
\draw (a)--(b)--(c)--(a);
\node[fill=none] at (0.5,-0.7) {$K_3$};
\end{scope}

\begin{scope}[xshift=8cm]
\node (a) at (0,0) {};
\node (b) at (1,0) {};
\node (c) at (1,1) {};
\node (d) at (0,1) {};
\draw (a)--(b)--(c)--(d)--(a);
\draw (a)--(c);
\draw (b)--(d);
\node[fill=none] at (0.5,-0.7) {$K_4$};
\end{scope}

\begin{scope}[xshift=11.5cm]
\node (a) at (0,0) {};
\node (b) at (1,0) {};
\node (c) at (1.3,1) {};
\node (d) at (0.5,1.6) {};
\node (e) at (-0.3,1) {};
\draw (a)--(b) (a)--(c) (a)--(d) (a)--(e);
\draw (b)--(c) (b)--(d) (b)--(e);
\draw (c)--(d) (c)--(e);
\draw (d)--(e);
\node[fill=none] at (0.5,-0.7) {$K_5$};
\end{scope}
\end{tikzpicture}
\caption{Grafos completos $K_1$ a $K_5$.}
\label{fig:completos}
\end{figure}

\begin{definicion}
Dado un grafo $G = (N,A)$, se define su \textbf{grafo complementario} $\overline{G}$ como aquel que tiene el mismo conjunto de nodos $N$ y cuyas conexiones son exactamente aquellas que no están presentes en $G$.

Así, para un grafo orientado:
\[
\overline{G} = (N, N \times N - A)
\]
mientras que para un grafo no orientado:
\[
\overline{G} = (N, \mathcal{P}_2(N) - A)
\]
\end{definicion}

\begin{observacion}
Nótese que el conjunto universal de referencia para las conexiones varía según el tipo de grafo:
\begin{itemize}
  \item En el caso \textbf{no orientado}, dicho universo corresponde al conjunto de todas las aristas posibles, es decir, las aristas de un grafo completo $K_n$.
  \item En el caso \textbf{orientado}, el universo estaría formado por todos los arcos posibles, incluyendo los lazos (es decir, $N \times N$); a este grafo se le denomina \textbf{grafo lleno}.
\end{itemize}

En la práctica es común excluir los lazos en el caso orientado. Bajo este convenio, el universo de referencia para la definición del complemento de un dígrafo es el \textbf{grafo completo simétrico} (aquel que contiene exactamente un arco en cada sentido entre todo par de nodos distintos). En el caso no orientado, la definición también excluye explícitamente los lazos, al tomar como universo $\mathcal{P}_2(N)$ en lugar de $\mathcal{P}(N)$.
\end{observacion}

\section{Trayectoria de un grafo}

\begin{definicion}
Un \textbf{recorrido} o \textbf{itinerario} o \textbf{cadena}, es una familia de conexiones sucesivamente adyacentes.
\end{definicion}

\begin{observacion}
\begin{itemize}
  \item Una cadena será \textbf{cerrada} si la última conexión de la sucesión fuera adyacente a la primera, y \textbf{abierta} en caso contrario. (Una cadena abierta puede contener subtrayectorias cerradas.)
  \item La notación es hecha indicando la sucesión a través de los nodos, las conexiones, o de nodos y conexiones alternados, o apenas los nodos inicial y final, cuando eso fuera suficiente. Por ejemplo, la cadena $\Pi_a$ puede ser expresada de las siguientes formas:
  \begin{align*}
  \Pi_a &= (1,2,3,4,5,3) \\
  \Pi_a &= ((1,2),(2,3),(3,4),(4,5),(5,3)) \\
  \Pi_a &= (1,(1,2),2,(2,3),3,(3,4),4,(4,5),5,(5,3),3) \\
  \Pi_a &= \Pi_{1,3}
  \end{align*}
  \item En un grafo no valorado, el \textbf{tamaño} de una cadena es el número de conexiones por ella utilizadas (contando las repeticiones). Para el caso de grafos valorados, será necesaria la generalización de distancias.
  \item Una cadena es \textbf{simple} si no repite conexiones, y es \textbf{elemental} si no repite nodos. La segunda noción es más restrictiva: toda cadena elemental es simple, pero lo recíproco no es verdadero.
  \item Un \textbf{ciclo} es una cadena simple y cerrada.
  \item Se denomina la \textbf{cintura} $g(G)$ de un grafo al tamaño del menor ciclo que este presente.
  \item Un \textbf{camino} es una cadena en un grafo orientado, en el cual la orientación de los arcos es siempre la misma, a partir del nodo inicial.
  \item Un \textbf{circuito} es un camino simple y cerrado en un grafo orientado.
\end{itemize}
\end{observacion}

\begin{definicion}
\begin{itemize}
  \item Una cadena (abierta o cerrada) se denomina \textbf{euleriana} si recorre cada conexión del grafo exactamente una vez; y se denomina \textbf{hamiltoniana} si visita cada nodo del grafo exactamente una vez.
  \item Si se permite la repetición, la cadena se denomina \textbf{pre-euleriana} o \textbf{pre-hamiltoniana}, según corresponda.
  \item Un grafo se dice \textbf{euleriano} (o \textbf{hamiltoniano}) si posee una cadena cerrada euleriana (o hamiltoniana, respectivamente).
\end{itemize}
\end{definicion}

\begin{ejemplo}[Cadena euleriana y grafo euleriano]
Considere el siguiente grafo no dirigido con nodos $A, B, C, D$ y ocho aristas entre ellos. Una cadena euleriana para este grafo es:
\[
A \to B \to D \to A \to C \to D \to B \to C \to A
\]
Esta cadena recorre cada una de las 8 aristas exactamente una vez y regresa al nodo inicial, por lo que es un \textbf{ciclo euleriano}. En consecuencia, el grafo es euleriano.
\end{ejemplo}

\begin{ejemplo}[Cadena hamiltoniana y grafo hamiltoniano]
Considere el siguiente grafo con nodos $A, B, C, D, E$. Una cadena hamiltoniana para este grafo es:
\[
A \to B \to C \to D \to E \to A
\]
Esta cadena visita cada uno de los 5 nodos exactamente una vez y regresa al nodo inicial, por lo que es un \textbf{ciclo hamiltoniano}. En consecuencia, el grafo es hamiltoniano.
\end{ejemplo}

\section{Grafos conexos -- Componentes f-conexas}

El concepto de conexidad está relacionado a la posibilidad de llegar de un nodo a otro en un grafo a través de las conexiones existentes. Esto es, el estado de conexión de un grafo adquiere aspectos diferentes conforme el grafo sea orientado o no. La idea de llegar está relacionada con la definición de alcanzable, la cual puede ser usada con interés, especialmente en grafos orientados.

\begin{definicion}
Un grafo (orientado o no) es \textbf{conexo} si para cualquier pareja de nodos del grafo se puede llegar al otro nodo partiendo de cualquiera de ellos.
\end{definicion}

\begin{definicion}
Un dígrafo en el cual todo par de nodos es unido por al menos una cadena se denomina \textbf{simplemente conexo} o \textbf{s-conexo}.
\end{definicion}

\begin{definicion}
Un dígrafo en el cual, en todo par de nodos, al menos uno es alcanzable a partir del otro, se denomina \textbf{semi-fuertemente conexo} o \textbf{sf-conexo}.
\end{definicion}

\begin{definicion}
Un dígrafo en el cual toda pareja de nodos son mutuamente alcanzables se denomina \textbf{fuertemente conexo} o \textbf{f-conexo}.
\end{definicion}

\begin{observacion}
\begin{enumerate}
  \item[a)] De las definiciones anteriores se desprende que todo dígrafo fuertemente conexo (f-conexo) es también semi-fuertemente conexo (sf-conexo) y, por tanto, simplemente conexo (s-conexo). Análogamente, todo dígrafo semi-fuertemente conexo es también simplemente conexo. Esta jerarquía permite establecer la siguiente clasificación, que constituye una partición del conjunto de todos los dígrafos:
  \begin{itemize}
    \item $C_3$ = dígrafos que son f-conexos
    \item $C_2$ = dígrafos que son sf-conexos y que no son f-conexos
    \item $C_1$ = dígrafos que son s-conexos y que no son sf-conexos
    \item $C_0$ = dígrafos que no son s-conexos
  \end{itemize}
  \item[b)] Las propiedades de los dígrafos de las categorías $C_1$ y $C_2$ ameritan una discusión en detalle por el interés en la presencia, en esos casos, de nodos privilegedos. La discusión de $C_3$ está asociada con el concepto de conexidad.
  \item[c)] Para un dígrafo $G$ y su complemento $\overline{G}$, se tiene:
  \begin{align*}
  G \in C_0 &\;\Rightarrow\; \overline{G} \in C_3 \\
  G \in C_1 \;\;\text{o}\;\; G \in C_2 \\
  G \in C_3 &\;\Rightarrow\; \overline{G} \notin C_i, \quad i = 0,1,2,3
  \end{align*}
\end{enumerate}
\end{observacion}

\begin{observacion}
La f-conexidad se caracteriza por la alcanzabilidad recíproca entre nodos, lo que convierte a la relación de alcanzabilidad en una relación simétrica. Además, todo nodo es alcanzable desde sí mismo (reflexividad), y si $z$ es alcanzable desde $y$, e $y$ es alcanzable desde $x$, entonces $z$ es alcanzable desde $x$ (transitividad). Por tanto, la relación de alcanzabilidad en un dígrafo $G = (N, A)$ es una relación de equivalencia sobre $N$, y en consecuencia induce una partición $S = \{S_1, S_2, \ldots, S_\gamma\}$ del conjunto de nodos. Los subgrafos inducidos por cada $S_i$ reciben el nombre de \textbf{componentes f-conexas maximales} o, simplemente, \textbf{componentes f-conexas}. Si el dígrafo es f-conexo, la partición consta de un único elemento: $S = \{N\}$.

La relación de sf-conexidad no es reflexiva; por tanto, las componentes sf-conexas no corresponden a una partición del conjunto de nodos $N$.
\end{observacion}

\section{Algunos resultados}

\begin{teorema}[Lema del apretón de manos]
Si $G = (N, A)$ es un multigrafo entonces
\[
\sum_{v \in N} \delta(v) = 2|A|
\]
\end{teorema}

\begin{proof}
Al considerar cada arista $\{a,b\}$ del grafo $G$ encontramos que la arista contribuye con una unidad a $\delta(a)$ y a $\delta(b)$, y en consecuencia, con dos unidades a $\sum_{v \in N} \delta(v)$. Así,
\[
\sum_{v \in N} \delta(v) = 2|A|. \qedhere
\]
\end{proof}

\begin{corolario}
Para cualquier multigrafo, el número de nodos de grado impar debe ser par.
\[
\sum_{\text{Par}} \delta(v) + \sum_{\text{Impar}} \delta(v) = 2|A|
\]
\end{corolario}

\begin{ejemplo}
¿Es posible tener un grafo 4-regular con 10 aristas?

\textbf{Solución.} De acuerdo con el teorema del apretón de manos, se tiene que:
\[
\sum_{v \in N} \delta(v) = 2|A|
\]
Si el grafo es 4-regular, entonces $\delta(v) = 4$ para todo $v \in N$. Por lo tanto:
\[
\sum_{v \in N} \delta(v) = 4|N|
\]
Sustituyendo en la fórmula:
\[
4|N| = 2(10) = 20 \quad\Longrightarrow\quad |N| = \frac{20}{4} = 5
\]
Por lo tanto, el grafo solicitado debe tener 5 nodos. Un ejemplo de tal grafo es el grafo completo $K_5$, que tiene 5 nodos y cada nodo tiene grado 4. El número de aristas de $K_5$ es:
\[
|A| = \binom{5}{2} = \frac{5 \cdot 4}{2} = 10
\]
Efectivamente, $K_5$ es un grafo 4-regular con 10 aristas.

\textbf{Respuesta:} Sí, es posible. El grafo $K_5$ cumple con las condiciones requeridas.
\end{ejemplo}

\begin{teorema}
En un dígrafo $G = (N,A)$, todo nodo del dígrafo se encuentra exactamente en una componente fuerte.
\end{teorema}

\begin{proof}
Las componentes f-conexas son subgrafos definidos en algún elemento de la partición definida por la relación de alcanzabilidad en $N$, es decir, que todo nodo del dígrafo se encuentra exactamente en una componente f-conexa.
\end{proof}

\begin{teorema}
Sea $G = (N,A)$ un grafo o multigrafo no dirigido sin nodos aislados. Entonces $G$ tiene un ciclo euleriano si y solo si $G$ es conexo y todo nodo de $G$ tiene grado par.
\end{teorema}

\begin{proof}
$(\Rightarrow)$ Sea $G$ euleriano. Entonces existe un ciclo que pasa por todo nodo $v \in N$; entonces, para un par de nodos cualesquiera $x$ e $y$, existe una cadena simple que puede llegar de $x$ a $y$; por tanto, $G$ es conexo. Por otro lado, al intentar recorrer el ciclo euleriano podremos escoger un nodo y de ahí proseguir atravesando los demás nodos, borrando las aristas utilizadas; al atravesar un nodo $x$, su grado $\delta(x)$ disminuirá en dos unidades; luego, todos los nodos intermediarios en la cadena deberán tener grado par, o será imposible anular todos sus grados al final de este proceso. El nodo inicial también deberá tener grado par, en vista que el uso de una de sus aristas adyacentes al inicio lo dejará con grado inicialmente impar, lo que permitirá su anulación al final del proceso.

$(\Leftarrow)$ Sea $G$ conexo, con todos sus nodos de grado par. Procederemos por inducción en el número de aristas. Si el número de aristas de $G$ es 1 o 2, entonces $G$ debe ser un ciclo, que evidentemente son eulerianos.

Ahora supongamos que $G$ es euleriano para cuando hay menos de $n$ aristas. Luego, si $G$ tiene $n$ aristas, seleccionamos un nodo $c$ en $G$ como punto inicial para construir un ciclo euleriano. Como $G$ es conexo y cada nodo tiene grado par, existe al menos un ciclo $\ell$ que contenga a $c$. Si $\ell$ tiene todas las aristas de $G$, hemos terminado. Si no, eliminamos las aristas del ciclo $\ell$ en $G$, asegurándonos de eliminar cualquier nodo que haya quedado aislado. El subgrafo restante $K$ tiene todos los nodos de grado par. Si $K$ es conexo, hemos terminado. Si $K$ no es conexo, cada componente de $K$ es conexa y tendrá un ciclo euleriano.

Además, cada uno de estos ciclos eulerianos tiene un nodo que está en $\ell$. En consecuencia, podemos partir de $c$ y recorrer $\ell$ hasta llegar al nodo $c_1$ que está en el ciclo euleriano de una componente de $K$; entonces se recorre este ciclo euleriano y, al regresar a $c_1$, continuamos en $\ell$ hasta llegar a un nodo $c_2$ que está en el ciclo euleriano de la otra componente de $K$. Como el grafo $G$ es finito, podemos continuar este proceso hasta construir un ciclo euleriano para $G$.
\end{proof}

\begin{corolario}
Si $G$ es un grafo o multigrafo no dirigido sin nodos aislados, entonces podemos construir una cadena abierta euleriana en $G$ si y solo si $G$ es conexo y tiene exactamente dos nodos de grado impar.
\end{corolario}

\begin{proof}
Si $G$ es conexo y $x$ e $y$ son los nodos de $G$ de grado impar, añadimos una arista $\{x,y\}$ a $G$. Ahora tenemos un grafo $G$ conexo tal que todos sus nodos son de grado par. Por lo tanto, $G$ tiene un ciclo euleriano $C$; cuando eliminamos la arista $\{x,y\}$ de $C$, obtenemos una cadena abierta euleriana de $x$ a $y$ en $G$. Lo recíproco se deja como tarea.
\end{proof}

\begin{ejemplo}[Aplicación del Teorema de Euler: el problema de los puentes de Königsberg]
Los habitantes se preguntaban si era posible realizar un recorrido que cruzara cada puente exactamente una vez y regresara al punto de partida. En términos de teoría de grafos, la pregunta equivale a determinar si este multigrafo (Figura~\ref{fig:konigsberg}) tiene un \textbf{ciclo euleriano}.

Aplicando el Teorema de Euler, debemos verificar dos condiciones:

\begin{enumerate}
  \item \textbf{El grafo debe ser conexo.} En efecto, el multigrafo que modela la ciudad es conexo, ya que es posible ir de cualquier zona a cualquier otra a través de los puentes.
  \item \textbf{Todos los nodos deben tener grado par.} Calculemos los grados de cada nodo:
  \[
  \delta(a) = 5, \quad \delta(b) = 3, \quad \delta(c) = 3, \quad \delta(d) = 3
  \]
\end{enumerate}

Observamos que todos los nodos tienen grado \textbf{impar}. Por lo tanto, el multigrafo \textbf{no satisface} la segunda condición del teorema. En consecuencia, no existe un ciclo euleriano que recorra cada puente exactamente una vez y regrese al punto de inicio.

Euler demostró que, para que tal recorrido sea posible, todos los nodos deben tener grado par, condición que no se cumple en este caso. Así quedó establecido que el paseo soñado por los habitantes de Königsberg era imposible, dando origen a la teoría de grafos.

\textbf{Nota:} si el objetivo hubiera sido un recorrido que cruzara cada puente exactamente una vez sin necesidad de regresar al punto de partida (camino euleriano abierto), el teorema establece que el grafo debe tener exactamente dos nodos de grado impar. En este caso, hay cuatro nodos de grado impar, por lo que tampoco existe un camino euleriano abierto.
\end{ejemplo}

\begin{corolario}
Sea $G = (N,A)$ un grafo o multigrafo dirigido sin nodos aislados. El grafo $G$ tiene un circuito euleriano dirigido si y solo si $G$ es conexo y $\delta^+(x) = \delta^-(x)$ para todo $x \in N$.
\end{corolario}

\begin{ejemplo}[Aplicación de circuitos eulerianos: generación de secuencias binarias]
Se muestra la superficie de un disco de rotación dividido en ocho sectores de igual área, con material conductor y no conductor colocado en el disco.

Cuando las tres terminales hacen contacto con tres sectores dados, el material no conductor produce un flujo de corriente y aparece un 1 en la pantalla de un dispositivo digital. Para los sectores con material conductor, se produce un flujo de corriente y aparece un 0 en la pantalla. Si el disco se rota 45 grados (en sentido de las manecillas del reloj), en la pantalla se leería 110 (de arriba a abajo). Así, podemos obtener al menos dos (100 y 110) de las ocho representaciones binarias del 0 al 7.

La pregunta es: ¿podemos representar los ocho números binarios al rotar el disco? Para responder esta pregunta, se construye un dígrafo $G = (N,A)$ donde
\[
N = \{00, 01, 10, 11\}
\]
y $A$ se construye como sigue: si $b_1b_2, b_2b_3 \in N$, entonces $(b_1b_2, b_2b_3) \in A$. Esto produce un dígrafo con lazos en $00$ y $11$, y arcos entre $00-01$, $01-10$ (doble), $01-11$, $10-00$, $10-11$.

Se observa que este dígrafo es conexo y que, para todo $x \in N$, $\delta^+(x) = \delta^-(x)$. Por el corolario anterior, se tiene un circuito euleriano:
\[
10 \to 00 \to 00 \to 01 \to 10 \to 01 \to 11 \to 11 \to 10
\]

Al etiquetar cada arco $a = (x,y)$ con la sucesión de tres bits $b_1b_2b_3$, donde $x = b_1b_2$ e $y = b_2b_3$, se determinan las ocho sucesiones de tres bits distintas. Así mismo, cualesquiera dos etiquetas de arcos consecutivos en el circuito euleriano son de la forma $b_1b_2b_3$ y $b_2b_3b_4$.

Partiendo del arco con etiqueta $100$, a fin de obtener la siguiente etiqueta ($000$), concatenamos el último bit de $000$, es decir $0$, a la cadena $100$. Entonces, la cadena resultante $1000$ proporciona $100$ ($100|0$) y $000$ ($1|000$). La siguiente etiqueta de arco es $001$, por lo que concatenamos el $1$ (el último bit de $001$) y obtenemos $10001$, lo cual proporciona las tres sucesiones de tres bits distintas: $100$ ($100|01$), $000$ ($1|000|1$) y $001$ ($10|001$). Continuando de esta forma llegamos a la sucesión de ocho bits:
\[
10001011
\]
(donde el último 1 se envuelve), y estos ocho bits se acomodan en los sectores del disco giratorio.

A partir de esta configuración, al rotar el disco, se obtienen las ocho sucesiones de tres bits:
\[
100,\ 110,\ 111,\ 011,\ 101,\ 010,\ 001,\ 000
\]
que corresponden a todas las representaciones binarias del 0 al 7.

Así, mediante el uso de un circuito euleriano en el dígrafo construido, se logra representar todos los números binarios del 0 al 7 al rotar el disco.
\end{ejemplo}

\subsection{Algoritmo de Fleury}

El algoritmo de Fleury obtiene un ciclo de Euler para un grafo conexo sin nodos de grado impar.

Una arista $\{x,y\}$ es un \textbf{puente} en un grafo conexo $G$ si al eliminar $\{x,y\}$ se crea un grafo disconexo.

Sea $G = (N,A)$ un grafo conexo con todos sus nodos de grado par.

\begin{algorithm}[htbp]
\caption{Algoritmo de Fleury}
\begin{algorithmic}[1]
\Require Un grafo conexo $G = (N,A)$ con todos sus nodos de grado par
\Ensure Un ciclo euleriano $\pi$
\State Elegir un nodo $v \in N$ como nodo inicial para el ciclo
\State $\pi \gets (v)$
\While{existan aristas en $A$}
  \State Suponer que ya se ha construido $\pi = (v, u, \ldots, w)$
  \If{en $w$ solo existe una arista $\{w,z\}$}
    \State Extender $\pi$ a $\pi = (v, u, \ldots, w, z)$
    \State Eliminar $\{w,z\}$ de $A$ y eliminar $w$ de $N$ si queda aislado
  \Else
    \State Elegir una arista $\{w,z\}$ que no sea un puente
    \State Extender $\pi$ a $\pi = (v, u, \ldots, w, z)$
    \State Eliminar $\{w,z\}$ de $A$
  \EndIf
\EndWhile
\State \Return $\pi$
\end{algorithmic}
\end{algorithm}

\subsection{Caminos Hamiltonianos}

Con respecto a las cadenas hamiltonianas surgen preguntas análogas a las correspondientes eulerianas. ¿Es posible determinar si existe una cadena hamiltoniana? Y si existe, ¿hay una manera eficiente de determinarla? Puede parecer sorprendente, pero la primera pregunta no ha sido resuelta de manera completa hasta la fecha. Sin embargo, daremos algunos resultados.

Es claro que los lazos y las aristas múltiples no son muy útiles para determinar las cadenas hamiltonianas, ya que entre dos nodos solo puede utilizarse una arista. También es fácil verificar que si un grafo $G$ de $n$ nodos tiene un ciclo hamiltoniano, entonces $G$ debe tener al menos $n$ aristas.

\begin{teorema}
Sea $G = (N,A)$ un grafo sin lazos, con $|N| = n \geq 2$. Si $\delta(x) + \delta(y) \geq n-1$ para todos $x, y \in N$, $x \neq y$, entonces $G$ tiene una cadena abierta hamiltoniana.
\end{teorema}

\begin{corolario}
Sea $G = (N,A)$ un grafo sin lazos, con $|N| = n \geq 2$. Si $\delta(x) \geq \dfrac{n-1}{2}$, $\forall x \in N$, entonces $G$ tiene una cadena abierta hamiltoniana.
\end{corolario}

\begin{teorema}
Sea $G = (N,A)$ un grafo sin lazos, con $|N| = n \geq 3$. Si $\delta(x) + \delta(y) \geq n$ para todos $x, y \in N$ no adyacentes, entonces $G$ tiene un ciclo hamiltoniano.
\end{teorema}

\begin{corolario}
Si $G = (N,A)$ es un grafo sin lazos, con $|N| = n \geq 3$ y si $\delta(x) \geq \dfrac{n}{2}$ para todo $x \in N$, entonces $G$ tiene un ciclo hamiltoniano.
\end{corolario}

\begin{corolario}
Si $G = (N,A)$ es un grafo sin lazos, con $|N| = n \geq 3$ y $|A| \geq \dfrac{1}{2}(n^2 - 3n + 6)$, entonces $G$ tiene un ciclo hamiltoniano.
\end{corolario}

Los recíprocos de los resultados anteriores no son válidos; es decir, las condiciones dadas son suficientes, pero no necesarias, para la conclusión.

El problema que se ha considerado tiene algunas variantes importantes. En un caso, las aristas pueden tener valores que representen una distancia, un costo o algo similar. Entonces el problema es determinar un ciclo hamiltoniano tal que la suma total de valores en la cadena sea mínima. Por ejemplo, los nodos podrían representar ciudades; las aristas, líneas de transporte; y el valor de una arista, el costo del viaje. Esta versión del problema se conoce como el \textbf{problema del agente viajero}.

\section{Un modelo de administración de proyectos}

En la década de 1950 aparecieron dos enfoques básicos para la planificación temporal, planteamiento y control de proyectos: PERT (\emph{Program Evaluation and Review Technique} o Técnica de Evaluación y Revisión de Programas) y CPM (\emph{Critical Path Method} o Método del camino crítico). Aunque en un principio existían diferencias fundamentales entre los dos métodos, el subsiguiente desarrollo de estos métodos ha disipado esas diferencias. Consiguientemente, en lo que sigue no se distingue entre PERT y CPM.

La construcción de una red de planificación temporal está limitada por las reglas siguientes:

\begin{itemize}
  \item Toda actividad debe ser representada por un único arco (esto es, las actividades son únicas).
  \item Como máximo, una actividad puede salir de un nodo y terminar en otro (esto es, el grafo es simple).
\end{itemize}

La segunda regla obliga al uso de \textbf{actividades mudas}. Una actividad muda no requiere tiempo alguno para su terminación, pero se necesita porque el grafo debe ser simple, y es necesario modelar ciertas relaciones lógicas.

El objetivo principal de la administración de proyectos es determinar el camino crítico y producir una planificación temporal que muestre los instantes de comienzo y finalización para todas las actividades de la red. Para la determinación de estas actividades temporales resultan útiles los siguientes conceptos:

\begin{definicion}
El \textbf{instante mínimo de comienzo}, $TE_j$, para cada suceso $j$, es el primer momento en que puede producirse ese suceso de tal manera que ya hayan finalizado todas las actividades entrantes de ese suceso.
\end{definicion}

El valor $TE_j$ de cada suceso se genera efectuando una pasada hacia adelante en la red, desde la fuente hasta el sumidero. En términos de la red, esto corresponde a la asignación de valores temporales a cada nodo (o suceso), de tal forma que el valor que se asigna a cada suceso es la cantidad de tiempo necesaria para completar las actividades presentes a lo largo del camino más largo que lleve a ese suceso.

Por definición, a la fuente se le asocia el valor cero. En general, el instante mínimo para un suceso está dado por una longitud máxima desde la fuente hasta el suceso.
\[
TE_1 = 0, \qquad TE_j = \max\{TE_i + t_{ij}\}, \quad j \neq 1
\]
para todas las posibles actividades que entren en el suceso $j$.

\begin{definicion}
El \textbf{tiempo máximo de finalización}, $TL_i$, es el último momento en que puede producirse un suceso sin dar lugar a un retraso del proyecto.
\end{definicion}

El valor de $TL_i$ para cada suceso se calcula efectuando una pasada hacia atrás desde el sumidero hacia la fuente.

El tiempo máximo de finalización asociado a cada suceso es el último momento en que puede finalizar una actividad sin dar lugar a un retraso en la fecha mínima de finalización del proyecto (esto es, son los tiempos finales de terminación asociados a aquellas actividades que no dan lugar a que aumente el valor $TE$ del nodo sumidero).

En general, el tiempo máximo de finalización para cualquier suceso se puede calcular de la siguiente manera:
\[
TL_s = TE_s, \qquad TL_i = \min\{TL_j - t_{ij}\}, \quad i \neq s
\]
donde $s$ es el sumidero. Es preciso tomar el mínimo porque al satisfacer la fecha final más temprana también se satisfacen las posteriores.

\section{Ejercicios}

\begin{ejercicio}
Los turistas Jenssen, Leuzinger, Dufour y Medeiros se encuentran en un bar de París y comienzan a conversar. Las lenguas disponibles son el inglés, el francés, el portugués y el alemán; Jenssen habla todas, Leuzinger no habla apenas el portugués, Dufour habla francés y alemán y Medeiros habla inglés y portugués.
\begin{enumerate}
  \item[a)] Represente por medio de un grafo $G = (N, A)$ todas las posibilidades de que uno de ellos dirija la palabra a otro, siendo comprendido. Defina $N$ y $A$.
  \item[b)] Represente por medio de un hipergrafo $H = (N, A')$ las capacidades lingüísticas del grupo. ¿Cuál es el significado de las intersecciones $a_i \cap a_j$ donde los $a_k \in A'$?
\end{enumerate}
\end{ejercicio}

\begin{ejercicio}
La figura muestra un grafo no dirigido que representa una sección de unos grandes almacenes. Los nodos indican el lugar donde se localizan las cajas; las aristas indican los pasillos que hay entre ellas. Los almacenes necesitan instalar un sistema de seguridad que consiste en colocar guardias (vestidos de civil) en ciertas cajas, de manera que cada cajero tenga un guardia en su lugar o que haya un solo pasillo entre una caja con guardia y el cajero. ¿Cuál es el número mínimo de guardias necesarios?
\end{ejercicio}

\begin{ejercicio}
Mostrar que los grafos dados son isomorfos.
\end{ejercicio}

\begin{ejercicio}
Mostrar que los dígrafos dados son isomorfos.
\end{ejercicio}

\begin{ejercicio}
Mostrar que los dígrafos dados no son isomorfos.
\end{ejercicio}

\begin{ejercicio}
Cada uno de los siguientes multigrafos surge en el análisis de un conjunto de cuatro bloques para el juego de ``locura instantánea''. Determine en cada caso si es posible resolver el acertijo.
\end{ejercicio}

\begin{ejercicio}
Diga si los dos grafos presentados son o no isomorfos. Si lo fueran, presente la función que establece el isomorfismo entre ellos; caso contrario, explique por qué.
\end{ejercicio}

\begin{ejercicio}
Escriba la matriz de adyacencia que representa a los siguientes grafos.
\end{ejercicio}

\begin{ejercicio}
Describa la matriz de adyacencias para $K_n$.
\end{ejercicio}

\begin{ejercicio}
Dada una matriz de adyacencias $A$ para un grafo simple $G$, describa la matriz de adyacencia para el complemento $\overline{G}$.
\end{ejercicio}

\begin{ejercicio}
Para el ejercicio 8, diseñe la representación en la forma de lista de adyacencias para el grafo indicado.
\end{ejercicio}

\begin{ejercicio}
Diseñe una representación en forma de matriz de incidencia para el grafo valorado de la figura.
\end{ejercicio}

\begin{ejercicio}
¿Cuál es el número total de aristas en $K_n$? Justifique su respuesta.
\end{ejercicio}

\begin{ejercicio}
¿Cuál es el grado de cada nodo de $K_n$?
\end{ejercicio}

\begin{ejercicio}
Si $G_1$ y $G_2$ son grafos no dirigidos (sin lazos), demuestre que $G_1, G_2$ son isomorfos si y solo si $\overline{G_1}$ y $\overline{G_2}$ son isomorfos. Determine si los siguientes grafos son isomorfos.
\end{ejercicio}

\begin{ejercicio}
Sea un grafo no dirigido con $n$ nodos. Si $G$ es isomorfo a su propio complemento $\overline{G}$ (este grafo se conoce como \textbf{autocomplementario}), ¿cuántas aristas debe tener $G$? Encuentre un ejemplo de un grafo autocomplementario con cuatro nodos y otro con cinco nodos. Si $G$ es un grafo autocomplementario con $n$ nodos, donde $n > 1$, demuestre que $n = 4k$ o $n = 4k+1$, para alguna $k \in \Z^+$.
\end{ejercicio}

\begin{ejercicio}
Sea $G$ un ciclo con $n$ nodos. Demuestre que $G$ es autocomplementario si y solo si $n = 5$.
\end{ejercicio}

\begin{ejercicio}
Dar tres caminos elementales que vayan de 1 a 3 para el dígrafo dado. ¿Existe algún circuito en el dígrafo? ¿Es transitivo el dígrafo?
\end{ejercicio}

\begin{ejercicio}
Calcular todos los semigrados exterior e interior de los nodos del dígrafo dado. Dar todos los circuitos elementales de este grafo. Obtener un dígrafo que no presente ningún circuito eliminando una arista del dígrafo dado. Determinar todos los nodos que puedan alcanzar cualquier otro nodo del dígrafo.
\end{ejercicio}

\begin{ejercicio}
Para los dígrafos dados en los ejercicios 18 y 19, determinar si son fuertemente conexos, semi-fuertemente conexos o simplemente conexos.
\end{ejercicio}

\begin{ejercicio}
Demostrar que un dígrafo $G$ es fuertemente conexo si y solo si existe un circuito en $G$ que incluye a todos los nodos al menos una vez.
\end{ejercicio}

\begin{ejercicio}
Hallar las componentes f-conexas del dígrafo dado en el problema 19. Calcular también sus componentes.
\end{ejercicio}

\begin{ejercicio}
El dígrafo de abajo representa las respuestas a la pregunta: ``¿Cuáles son los colegas de quien Ud. más gusta?'' dadas por una promoción de primer semestre de estudiantes de la UNSA (un grafo como este es denominado \textbf{sociograma}). Use notación y nomenclatura convenientes para indicar la existencia de líder(es), amistades recíprocas, problemas de relacionamiento, influencias directas o indirectas y aislamiento.
\end{ejercicio}

\begin{ejercicio}
Un museo de arte ha ordenado la exposición que actualmente presenta en cinco salas, como muestra la figura. ¿Existe alguna forma de recorrer la exposición de modo que Ud. pase por cada puerta solo una vez? En este caso, trace el recorrido.
\end{ejercicio}

\begin{ejercicio}
En la puerta de una mansión histórica, usted recibe una copia del plano de la casa. ¿Es posible visitar cada cuarto de la casa pasando por cada puerta solo una vez? Explique el razonamiento.
\end{ejercicio}

\begin{ejercicio}
Trace $K_{2,3}$, $K_{2,4}$, $K_{3,3}$.
\end{ejercicio}

\begin{ejercicio}
Obtenga una fórmula para el número de aristas de $K_{m,n}$.
\end{ejercicio}

\begin{ejercicio}
¿Para qué valores de $m$ y $n$ tiene $K_{m,n}$ un circuito euleriano?
\end{ejercicio}

\begin{ejercicio}
Diseñe un juego de ``locura instantánea'' que tenga exactamente cuatro soluciones.
\end{ejercicio}

\begin{ejercicio}
En el siguiente grafo, los nodos representan ciudades y los números de las aristas señalan los costos de construcción de la carretera correspondiente. Determine un sistema de caminos de menor costo que conecte todas las ciudades. ¿Es esta una solución única?
\end{ejercicio}

\begin{ejercicio}
Suponga que un grafo tiene una matriz de adyacencia de la forma
\[
A =
\begin{bmatrix}
A' & 0 \\
0 & A''
\end{bmatrix}
\]
en donde los elementos fuera de las submatrices $A'$ y $A''$ son ceros. ¿Cómo debe ser el esquema del grafo?
\end{ejercicio}

\begin{ejercicio}
Demuestre que un camino que no es simple debe contener un circuito.
\end{ejercicio}

\begin{ejercicio}
Demuestre que si $G'$ es un subgrafo conexo de un grafo $G$, entonces $G'$ está contenido en un componente de $G$.
\end{ejercicio}

\begin{ejercicio}
Todo nodo y toda arista de un dígrafo están contenidos en una componente s-conexa y en solo una.
\end{ejercicio}

\begin{ejercicio}
Demostrar el Corolario 5.36.
\end{ejercicio}

\begin{ejercicio}
Determine los valores de $n$ para los que el grafo completo $K_n$ tiene ciclo euleriano. ¿Para cuáles $n$ tiene $K_n$ una cadena abierta euleriana pero no un ciclo euleriano?
\end{ejercicio}

\begin{ejercicio}
Diga si es posible trazar la figura sin levantar el lápiz. Explique su razonamiento.
\end{ejercicio}

\begin{ejercicio}
Sea $N = \{000, 001, 010, \ldots, 111\}$. Para cualquier sucesión de cuatro bits $b_1, b_2, b_3, b_4$, trace una arista del elemento $b_1b_2b_3$ al elemento $b_2b_3b_4$ en $N$.
\begin{enumerate}
  \item[a)] Trace el grafo $G = (N,A)$ descrito.
  \item[b)] Encuentre un ciclo euleriano para $G$.
  \item[c)] Distribuya ocho ceros y ocho unos de modo uniforme alrededor del borde de un disco que gira en el sentido de las manecillas del reloj, de modo que estos 16 bits formen una sucesión circular tal que las subsucesiones (consecutivas) de longitud 4 proporcionen las representaciones binarias de $0,1,2,3,\ldots,15$ en algún orden.
\end{enumerate}
\end{ejercicio}

\begin{ejercicio}
Determine $|N|$ para los siguientes grafos o multigrafos $G$:
\begin{enumerate}
  \item[a)] $G$ tiene nueve aristas y todos los nodos tienen grado 3.
  \item[b)] $G$ es regular con 15 aristas.
  \item[c)] $G$ tiene 10 aristas con dos nodos de grado 4 y los demás de grado 3.
\end{enumerate}
\end{ejercicio}

\begin{ejercicio}
Si $G = (N,A)$ es un grafo conexo con $|A| = 17$ y $\delta(x) \geq 3$ para todo $x \in N$, ¿cuál es el valor máximo para $|N|$?
\end{ejercicio}

\begin{ejercicio}
Sea $N = \{a,b,c,d,e,f\}$. Dibuje tres grafos sin lazos en los que $\delta(a) = 3$, $\delta(b) = \delta(c) = 2$ y $\delta(d) = \delta(e) = \delta(f) = 1$. ¿Cuántos de esos grafos son conexos?
\end{ejercicio}

\begin{ejercicio}
Dé un ejemplo de un grafo conexo tal que:
\begin{enumerate}
  \item[a)] No tenga ciclos eulerianos ni ciclos hamiltonianos.
  \item[b)] Tenga un ciclo euleriano pero no un ciclo hamiltoniano.
  \item[c)] Tenga un ciclo hamiltoniano pero no un ciclo euleriano.
  \item[d)] Tenga un ciclo hamiltoniano y un ciclo euleriano.
\end{enumerate}
\end{ejercicio}

\begin{ejercicio}
Determine si el grafo dado tiene un ciclo hamiltoniano, una cadena hamiltoniana pero no un ciclo hamiltoniano, o ninguno de los dos. Si el grafo tiene un ciclo hamiltoniano, indique el ciclo.
\end{ejercicio}
\chapter{Conexión y Recorridos}

\section{Trayectorias y Caminos}
Un camino en un grafo es una secuencia de vértices donde cada par consecutivo está conectado por una arista.

\subsection{Tipos de caminos}
\begin{itemize}
  \item \textbf{Camino simple}: No repite vértices.
  \item \textbf{Ciclo}: Camino que comienza y termina en el mismo vértice.
  \item \textbf{Trayectoria}: Camino sin aristas repetidas.
\end{itemize}

\section{Grafos Conexos}
Un grafo es conexo si existe un camino entre cualquier par de vértices.

\subsection{Componentes conexas}
Una componente conexa es un subgrafo maximal que es conexo.

\section{Recorridos en Grafos}

\subsection{Búsqueda en profundidad (DFS)}

\begin{algorithm}[H]
\caption{DFS (Búsqueda en Profundidad)}
\label{alg:dfs}
\begin{algorithmic}[1]
\Require Un grafo $G=(V,E)$ y un vértice inicial $s$
\Ensure Recorrido DFS
\State visitado$[v] \gets \text{falso}$ para todo $v \in V$
\Procedure{DFS}{$u$}
  \State visitado$[u] \gets \text{verdadero}$
  \For{cada $v \in \text{adyacentes}(u)$}
    \If{$\neg$visitado$[v]$}
      \State DFS$(v)$
    \EndIf
  \EndFor
\EndProcedure
\State DFS$(s)$
\end{algorithmic}
\end{algorithm}

\subsection{Búsqueda en anchura (BFS)}

\begin{algorithm}[H]
\caption{BFS (Búsqueda en Anchura)}
\label{alg:bfs}
\begin{algorithmic}[1]
\Require Un grafo $G=(V,E)$ y un vértice inicial $s$
\State visitado$[v] \gets \text{falso}$ para todo $v \in V$
\State visitado$[s] \gets \text{verdadero}$
\State cola $\gets [s]$
\While{cola $\neq \emptyset$}
  \State $u \gets \text{POPFRONT}(cola)$
  \For{cada $v \in \text{adyacentes}(u)$}
    \If{$\neg$visitado$[v]$}
      \State visitado$[v] \gets \text{verdadero}$
      \State cola $\gets$ cola $+$ $[v]$
    \EndIf
  \EndFor
\EndWhile
\end{algorithmic}
\end{algorithm}

\section{Caminos Eulerianos y Hamiltonianos}

\subsection{Caminos Eulerianos}
Un camino euleriano recorre todas las aristas exactamente una vez. Un grafo conexo tiene un ciclo euleriano si y solo si todos los vértices tienen grado par.

\subsection{Caminos Hamiltonianos}
Un camino hamiltoniano visita todos los vértices exactamente una vez. Teorema de Dirac: Si un grafo con $n \geq 3$ vértices tiene $\deg(v) \geq n/2$ para todo $v$, entonces tiene un ciclo hamiltoniano.

\section{Implementación en Haskell}

\begin{lstlisting}[caption={Algoritmos de recorrido en Haskell}, label=lst:grafos-haskell]
-- DFS
dfs :: Eq a => a -> [(a, [a])] -> [a]
dfs inicio grafo = dfsAux inicio []
  where
    dfsAux v visitados
      | v `elem` visitados = visitados
      | otherwise = foldl (\acc w -> dfsAux w acc) (v:visitados) (adyacentes v grafo)

-- BFS
bfs :: Eq a => a -> [(a, [a])] -> [a]
bfs inicio grafo = bfsAux [inicio] []
  where
    bfsAux [] visitados = visitados
    bfsAux (v:cola) visitados
      | v `elem` visitados = bfsAux cola visitados
      | otherwise = bfsAux (cola ++ adyacentes v grafo) (visitados ++ [v])

-- Verificar si es conexo
esConexo :: Eq a => [(a, [a])] -> Bool
esConexo grafo = all (`elem` reachable) (map fst grafo)
  where (inicio:_) = map fst grafo
        reachable = dfs inicio grafo

-- Verificar si tiene ciclo euleriano
cicloEuleriano :: Eq a => [(a, [a])] -> Bool
cicloEuleriano grafo = esConexo grafo && all even (map (length . snd) grafo)
\end{lstlisting}

\section{Ejercicios}
\begin{enumerate}
  \item Ejecutar DFS y BFS en un grafo de ejemplo.
  \item Demostrar que un grafo tiene un ciclo euleriano si todos los vértices tienen grado par.
  \item Implementar en Haskell un algoritmo que encuentre un camino euleriano.
\end{enumerate}
\chapter{Árboles}

\section{Definiciones Básicas}
Un \textbf{árbol} es un grafo conexo sin ciclos. Para un grafo $G$ con $n$ vértices son equivalentes:
\begin{enumerate}
  \item $G$ es un árbol.
  \item $G$ es conexo y tiene $n-1$ aristas.
  \item $G$ no tiene ciclos y tiene $n-1$ aristas.
  \item Entre cualquier par de vértices existe exactamente un camino.
\end{enumerate}

\section{Clasificación}

\subsection{Árboles binarios}
Un árbol binario es un árbol enraizado donde cada nodo tiene a lo sumo dos hijos. Propiedades:
\begin{itemize}
  \item Un árbol binario con $n$ nodos tiene $n-1$ aristas.
  \item La altura es al menos $\lceil \log_2(n+1) \rceil$.
  \item El nivel $i$ tiene a lo sumo $2^i$ nodos.
\end{itemize}

\section{Árboles Ponderados}

\subsection{Árboles generadores}
Un \textbf{árbol generador} de $G$ es un subgrafo que es árbol y contiene todos los vértices.

\subsection{Códigos de Huffman}

\begin{algorithm}[htbp]
\caption{Algoritmo de Huffman}
\begin{algorithmic}[1]
\Require Símbolos con frecuencias
\Ensure Árbol de codificación binario
\State Crear nodo hoja por símbolo con su frecuencia
\State Insertar todos los nodos en cola de prioridad
\While{cola tiene más de un elemento}
  \State Extraer los dos nodos de menor frecuencia
  \State Crear nuevo nodo con frecuencia suma
  \State Los nodos extraídos son hijos del nuevo
  \State Insertar el nuevo en la cola
\EndWhile
\State \Return raíz del árbol
\end{algorithmic}
\end{algorithm}

\section{Implementación en Haskell}

\begin{lstlisting}[caption={Árboles binarios en Haskell}]
data ArbolBinario a = Vacio | Nodo a (ArbolBinario a) (ArbolBinario a)
  deriving (Show, Eq)

preorden :: ArbolBinario a -> [a]
preorden Vacio = []
preorden (Nodo x izq der) = x : preorden izq ++ preorden der

inorden :: ArbolBinario a -> [a]
inorden Vacio = []
inorden (Nodo x izq der) = inorden izq ++ [x] ++ inorden der

postorden :: ArbolBinario a -> [a]
postorden Vacio = []
postorden (Nodo x izq der) = postorden izq ++ postorden der ++ [x]

altura :: ArbolBinario a -> Int
altura Vacio = 0
altura (Nodo _ izq der) = 1 + max (altura izq) (altura der)
\end{lstlisting}

\section{Ejercicios}
\begin{enumerate}
  \item Demostrar que un árbol con $n$ vértices tiene $n-1$ aristas.
  \item Implementar en Haskell un árbol de Huffman.
  \item ¿Cuál es la altura máxima de un árbol binario con $n$ nodos?
\end{enumerate}
\chapter{Algoritmos sobre Grafos}

\section{Camino más Corto}

\subsection{Algoritmo de Dijkstra}

\begin{algorithm}[htbp]
\caption{Algoritmo de Dijkstra}
\begin{algorithmic}[1]
\Require Grafo ponderado $G=(V,E)$, origen $s$
\Ensure Distancias mínimas desde $s$
\State dist$[v] \gets \infty$, dist$[s] \gets 0$
\State $Q \gets V$
\While{$Q \neq \emptyset$}
  \State $u \gets \arg\min_{v \in Q}$ dist$[v]$
  \State $Q \gets Q \setminus \{u\}$
  \For{cada $v$ adyacente a $u$}
    \If{dist$[v] >$ dist$[u] + w(u,v)$}
      \State dist$[v] \gets$ dist$[u] + w(u,v)$
    \EndIf
  \EndFor
\EndWhile
\State \Return dist
\end{algorithmic}
\end{algorithm}

\section{Árboles Generadores Mínimos}

\begin{algorithm}[htbp]
\caption{Algoritmo de Prim}
\begin{algorithmic}[1]
\Require Grafo ponderado conexo $G=(V,E)$
\Ensure Árbol generador mínimo $T$
\State $T \gets \emptyset$, $S \gets \{v_0\}$
\While{$S \neq V$}
  \State Encontrar $(u,v)$ de peso mínimo con $u \in S$, $v \notin S$
  \State $S \gets S \cup \{v\}$
  \State $T \gets T \cup \{(u,v)\}$
\EndWhile
\State \Return $T$
\end{algorithmic}
\end{algorithm}

\begin{algorithm}[htbp]
\caption{Algoritmo de Kruskal}
\begin{algorithmic}[1]
\Require Grafo ponderado conexo $G=(V,E)$
\Ensure Árbol generador mínimo $T$
\State $T \gets \emptyset$
\State Ordenar $E$ por peso ascendente
\For{cada $(u,v) \in E$}
  \If{$u$ y $v$ en componentes distintas}
    \State $T \gets T \cup \{(u,v)\}$
    \State Unir componentes de $u$ y $v$
  \EndIf
\EndFor
\State \Return $T$
\end{algorithmic}
\end{algorithm}

\section{Redes de Transporte}

\textbf{Teorema de Flujo Máximo - Corte Mínimo}: el flujo máximo en una red es igual a la capacidad del corte mínimo.

\begin{algorithm}[htbp]
\caption{Ford-Fulkerson}
\begin{algorithmic}[1]
\Require Red con fuente $s$, sumidero $t$, capacidades $c$
\Ensure Flujo máximo
\State flujo$[u][v] \gets 0$ para todo $u,v$
\While{existe camino de aumento $P$ de $s$ a $t$}
  \State Encontrar capacidad residual mínima en $P$
  \State Actualizar flujo en $P$
\EndWhile
\State \Return flujo
\end{algorithmic}
\end{algorithm}

\section{Ejercicios}
\begin{enumerate}
  \item Ejecutar Dijkstra en un grafo de ejemplo.
  \item Demostrar que Prim encuentra un árbol generador mínimo.
  \item Implementar Kruskal en Haskell.
\end{enumerate}

\part{Criptografía}
\chapter{Aritmética Modular}

\section{Introducción}
$a \equiv b \pmod{n}$ si $n$ divide a $a - b$.

Propiedades: reflexiva, simétrica, transitiva.

\section{Operaciones Modulares}
\[
a + b \equiv (a \bmod n) + (b \bmod n) \pmod{n},
\]
\[
a \cdot b \equiv (a \bmod n) \cdot (b \bmod n) \pmod{n},
\]
\[
a^b \bmod n.
\]

\section{Algoritmo de Euclides}

\begin{algorithm}[htbp]
\caption{Algoritmo de Euclides}
\begin{algorithmic}[1]
\Require $a \geq b > 0$
\Ensure $\gcd(a,b)$
\While{$b \neq 0$}
  \State $r \gets a \bmod b$
  \State $a \gets b$
  \State $b \gets r$
\EndWhile
\State \Return $a$
\end{algorithmic}
\end{algorithm}

El algoritmo extendido encuentra $x, y$ tales que $ax + by = \gcd(a,b)$.

\section{Inversos Modulares}
$a^{-1}$ existe si y solo si $\gcd(a,n) = 1$.

\section{Pequeño Teorema de Fermat}
Si $p$ es primo y $p \nmid a$:
\[
a^{p-1} \equiv 1 \pmod{p}.
\]

\section{Implementación en Haskell}

\begin{lstlisting}[caption={Aritmética modular en Haskell}]
mcd :: Integral a => a -> a -> a
mcd a 0 = abs a
mcd a b = mcd b (a `mod` b)

euclidesExtendido :: Integral a => a -> a -> (a, a, a)
euclidesExtendido a 0 = (a, 1, 0)
euclidesExtendido a b =
    let (g, x1, y1) = euclidesExtendido b (a `mod` b)
    in (g, y1, x1 - (a `div` b) * y1)

inversoModular :: Integral a => a -> a -> Maybe a
inversoModular a m =
    let (g, x, _) = euclidesExtendido a m
    in if g == 1 then Just (x `mod` m) else Nothing

potenciaMod :: Integer -> Integer -> Integer -> Integer
potenciaMod base 0 _ = 1
potenciaMod base exp mod
    | even exp = let p = potenciaMod base (exp `div` 2) mod
                 in (p*p) `mod` mod
    | otherwise = (base * potenciaMod base (exp - 1) mod) `mod` mod
\end{lstlisting}

\section{Ejercicios}
\begin{enumerate}
  \item Calcular el inverso modular de 3 módulo 11.
  \item Demostrar el Pequeño Teorema de Fermat.
  \item Implementar en Haskell una función que verifique si un número es primo usando el test de Fermat.
\end{enumerate}
\chapter{Criptografía Clásica}

\section{Introducción}
\begin{itemize}
  \item \textbf{Criptografía simétrica}: misma clave para cifrar y descifrar.
  \item \textbf{Criptografía asimétrica}: claves diferentes.
\end{itemize}

\section{Cifrado César}
Desplaza cada letra un número fijo de posiciones. Con desplazamiento 3:
\[
\text{HOLA} \to \text{KROD}.
\]
Solo tiene 25 desplazamientos: es vulnerable a fuerza bruta.

\section{Cifrado de Sustitución}
Reemplaza cada letra por otra según una tabla. El \textbf{análisis de frecuencia} permite romperlo.

\section{Cifrado de Vigenère}
Usa una palabra clave para desplazar cada letra de forma diferente. Con clave \texttt{CLAVE}:
\[
\text{HOLA} \to \text{JPRD}.
\]

\section{Implementación en Haskell}

\begin{lstlisting}[caption={Criptografía clásica en Haskell}]
import Data.Char

cesar :: Int -> String -> String
cesar n = map (cifrar n)
  where
    cifrar n c
      | isAlpha c = chr (ord 'A' + ((ord (toUpper c) - ord 'A' + n) `mod` 26))
      | otherwise = c

cesarFuerzaBruta :: String -> [String]
cesarFuerzaBruta texto = [cesar n texto | n <- [0..25]]

vigenere :: String -> String -> String
vigenere clave texto =
    zipWith cifrar (map toUpper texto) (cycle (map toUpper clave))
  where
    cifrar c k = chr (ord 'A' + ((ord c - ord 'A' + ord k - ord 'A') `mod` 26))
\end{lstlisting}

\section{Ejercicios}
\begin{enumerate}
  \item Cifrar y descifrar con César.
  \item Romper un cifrado César por fuerza bruta.
  \item Implementar Vigenère en Haskell.
\end{enumerate}
\chapter{Criptografía RSA}

\section{Introducción}
RSA es uno de los sistemas criptográficos más importantes. Fue desarrollado por Rivest, Shamir y Adleman en 1977.

El sistema RSA utiliza dos claves: una pública $(e,n)$ y una privada $(d,n)$.

\section{Generación de Claves}
El proceso de generación de claves RSA consiste en:
\begin{enumerate}
  \item Elegir dos números primos grandes $p$ y $q$.
  \item Calcular $n = p \cdot q$.
  \item Calcular $\varphi(n) = (p-1)(q-1)$.
  \item Elegir $e$ tal que $1 < e < \varphi(n)$ y $\gcd(e,\varphi(n)) = 1$.
  \item Calcular $d = e^{-1} \bmod \varphi(n)$.
\end{enumerate}

\section{Cifrado y Descifrado}

\subsection{Cifrado}
Para cifrar un mensaje $M$:
\[
C = M^e \bmod n.
\]

\subsection{Descifrado}
Para descifrar un mensaje $C$:
\[
M = C^d \bmod n.
\]

\subsection{Demostración de corrección}
La corrección de RSA se basa en el teorema de Euler:
\[
M^{ed} \equiv M \pmod{n}.
\]

\section{Seguridad de RSA}
La seguridad de RSA se basa en la dificultad de factorizar $n = p \cdot q$.

\section{Implementación en Haskell}

\begin{lstlisting}[caption={RSA en Haskell}, label=lst:rsa-haskell]
-- Algoritmo de Euclides extendido
euclidesExt :: Integer -> Integer -> (Integer, Integer, Integer)
euclidesExt a 0 = (a, 1, 0)
euclidesExt a b =
    let (g, x1, y1) = euclidesExt b (a `mod` b)
    in (g, y1, x1 - (a `div` b) * y1)

-- Inverso modular
inversoMod :: Integer -> Integer -> Maybe Integer
inversoMod a m =
    let (g, x, _) = euclidesExt a m
    in if g == 1 then Just (x `mod` m) else Nothing

-- Potenciación modular
potenciaMod :: Integer -> Integer -> Integer -> Integer
potenciaMod base 0 _ = 1
potenciaMod base exp mod
    | even exp = (potenciaMod base (exp `div` 2) mod) ^ 2 `mod` mod
    | otherwise = (base * potenciaMod base (exp - 1) mod) `mod` mod

-- Verificar si es primo (simple)
esPrimo :: Integer -> Bool
esPrimo n = n > 1 && all (\p -> n `mod` p /= 0) (takeWhile (<= floor (sqrt (fromIntegral n))) [2..])

-- Generar clave RSA
generarRSA :: Integer -> Integer -> Integer -> ((Integer, Integer), (Integer, Integer))
generarRSA p q e =
    let n = p * q
        phi = (p-1) * (q-1)
        d = case inversoMod e phi of
              Just d' -> d'
              Nothing -> error "e no es invertible módulo phi(n)"
    in ((e, n), (d, n))

-- Cifrar RSA
cifrarRSA :: Integer -> (Integer, Integer) -> Integer
cifrarRSA m (e, n) = potenciaMod m e n

-- Descifrar RSA
descifrarRSA :: Integer -> (Integer, Integer) -> Integer
descifrarRSA c (d, n) = potenciaMod c d n

-- Conversión texto-número
textoANumero :: String -> Integer
textoANumero = foldl (\acc c -> acc * 256 + fromIntegral (ord c)) 0

numeroATexto :: Integer -> String
numeroATexto 0 = ""
numeroATexto n = let (q, r) = n `divMod` 256 in numeroATexto q ++ [chr (fromIntegral r)]

-- Ejemplo
ejemploRSA :: IO ()
ejemploRSA = do
    let p = 61
        q = 53
        e = 17
        ((pubE, pubN), (privD, privN)) = generarRSA p q e
        mensaje = 123
        cifrado = cifrarRSA mensaje (pubE, pubN)
        descifrado = descifrarRSA cifrado (privD, privN)
    putStrLn $ "Clave pública: (" ++ show pubE ++ ", " ++ show pubN ++ ")"
    putStrLn $ "Clave privada: (" ++ show privD ++ ", " ++ show privN ++ ")"
    putStrLn $ "Mensaje original: " ++ show mensaje
    putStrLn $ "Mensaje cifrado: " ++ show cifrado
    putStrLn $ "Mensaje descifrado: " ++ show descifrado
\end{lstlisting}

\section{Ejercicios}
\begin{enumerate}
  \item Generar un par de claves RSA con primos pequeños.
  \item Cifrar y descifrar un mensaje.
  \item Implementar en Haskell una función que convierta texto a números para RSA.
\end{enumerate}

\part{Programación Funcional con Haskell}
\chapter{Introducción a Haskell}

\section{¿Qué es Haskell?}
Haskell es un lenguaje funcional puro, de propósito general. Se utiliza en verificación de hardware, telecomunicaciones, comercio electrónico, etc.

\subsection{Características principales}
\begin{itemize}
  \item \textbf{Transparencia referencial}: el resultado depende solo de los argumentos.
  \item \textbf{Funciones como ciudadanos de primera clase}: se pasan, devuelven y almacenan.
  \item \textbf{Evaluación perezosa}: se evalúa solo cuando se necesita.
  \item \textbf{Tipos fuerte y estático}: verificación en compilación.
  \item \textbf{Inferencia de tipos}.
  \item \textbf{Polimorfismo paramétrico y sobrecarga}.
\end{itemize}

\section{Tipos Básicos}

\subsection{Enteros}
\begin{table}[h]
\centering
\begin{tabular}{|l|p{5cm}|p{5cm}|}
\hline
\textbf{Característica} & \textbf{Int} & \textbf{Integer} \\
\hline
Tamaño & Fijo & Arbitrario \\
Rango & Limitado & Ilimitado \\
Rendimiento & Rápido & Lento \\
Riesgo & Desbordamiento & Ninguno \\
\hline
\end{tabular}
\caption{Comparación entre Int e Integer}
\end{table}

\subsection{Reales}
\texttt{Float} (7 dígitos), \texttt{Double} (15-16 dígitos).

\subsection{Bool, Char, String}
\texttt{Bool}: \texttt{True}, \texttt{False}. \texttt{Char}: \texttt{'a'}. \texttt{String} es \texttt{[Char]}.

\section{Definición de funciones}

\begin{lstlisting}
areaRectangulo :: Int -> Int -> Int
areaRectangulo b h = b * h

maximo :: Int -> Int -> Int
maximo x y | x <= y = y
           | otherwise = x

factorial :: Integer -> Integer
factorial 0 = 1
factorial n = n * factorial (n - 1)
\end{lstlisting}

\section{Listas}

\begin{lstlisting}
[1,2,3] :: [Int]
[1..10] -- [1,2,3,4,5,6,7,8,9,10]
[1,3..10] -- [1,3,5,7,9]

dobles :: [Int] -> [Int]
dobles xs = [2*x | x <- xs]

map :: (a -> b) -> [a] -> [b]
map (+1) [1,2,3] -- [2,3,4]
\end{lstlisting}

\section{Funciones de Orden Superior}

\begin{lstlisting}
dosVeces :: (a -> a) -> a -> a
dosVeces f x = f (f x)

(.) :: (b -> c) -> (a -> b) -> (a -> c)
(g . f) x = g (f x)
\end{lstlisting}

\section{Tuplas}

\begin{lstlisting}
(1, "hola") :: (Int, String)
fst (1, "hola") -- 1
snd (1, "hola") -- "hola"
\end{lstlisting}

\section{Ejercicios}
(Ver el texto completo.)
\chapter{Haskell para Combinatoria}

\section{Generación de Permutaciones}

\begin{lstlisting}[caption={Permutaciones en Haskell}]
permutaciones :: [a] -> [[a]]
permutaciones [] = [[]]
permutaciones xs = [ y:ys | y <- xs, ys <- permutaciones (eliminar y xs) ]
  where
    eliminar _ [] = []
    eliminar y (x:xs') | y == x = xs'
                       | otherwise = x : eliminar y xs'

combinaciones :: Int -> [a] -> [[a]]
combinaciones 0 _ = [[]]
combinaciones _ [] = []
combinaciones k (x:xs) = map (x:) (combinaciones (k-1) xs) ++ combinaciones k xs

subconjuntos :: [a] -> [[a]]
subconjuntos [] = [[]]
subconjuntos (x:xs) = subconjuntos xs ++ map (x:) (subconjuntos xs)
\end{lstlisting}

\section{Números Especiales}

\begin{lstlisting}
catalan :: Integer -> Integer
catalan n = product [n+2..2*n] `div` product [1..n]

stirling2 :: Integer -> Integer -> Integer
stirling2 0 0 = 1
stirling2 _ 0 = 0
stirling2 0 _ = 0
stirling2 n k = k * stirling2 (n-1) k + stirling2 (n-1) (k-1)
\end{lstlisting}
\chapter{Haskell para Grafos}

\section{Representación de Grafos}

\begin{lstlisting}[caption={Representación de grafos en Haskell}]
type Vertice = Int
type Grafo = [(Vertice, [Vertice])]
type GrafoPonderado = [(Vertice, [(Vertice, Int)])]

grafoDesdeAristas :: [Vertice] -> [(Vertice,Vertice)] -> Grafo
grafoDesdeAristas vs aristas =
    [(v, [w | (a,w) <- aristas, a == v] ++ [w | (w,a) <- aristas, a == v]) | v <- vs]

grafoDirigido :: [Vertice] -> [(Vertice,Vertice)] -> Grafo
grafoDirigido vs aristas =
    [(v, [w | (a,w) <- aristas, a == v]) | v <- vs]
\end{lstlisting}

\section{Algoritmos de Recorrido}

\begin{lstlisting}[caption={DFS y BFS en Haskell}]
dfs :: Vertice -> Grafo -> [Vertice]
dfs inicio grafo = dfsAux inicio []
  where
    dfsAux v vis
      | v `elem` vis = vis
      | otherwise = foldl (\acc w -> dfsAux w acc) (v:vis)
                    (maybe [] id (lookup v grafo))

bfs :: Vertice -> Grafo -> [Vertice]
bfs inicio grafo = bfsAux [inicio] []
  where
    bfsAux [] vis = vis
    bfsAux (v:cola) vis
      | v `elem` vis = bfsAux cola vis
      | otherwise = bfsAux (cola ++ maybe [] id (lookup v grafo)) (vis ++ [v])
\end{lstlisting}
\chapter{Haskell para Criptografía}

\section{Aritmética Modular}

\begin{lstlisting}[caption={Aritmética modular en Haskell}]
euclidesExt :: Integer -> Integer -> (Integer, Integer, Integer)
euclidesExt a 0 = (a, 1, 0)
euclidesExt a b =
    let (g, x1, y1) = euclidesExt b (a `mod` b)
    in (g, y1, x1 - (a `div` b) * y1)

inversoMod :: Integer -> Integer -> Maybe Integer
inversoMod a m =
    let (g, x, _) = euclidesExt a m
    in if g == 1 then Just (x `mod` m) else Nothing

potenciaMod :: Integer -> Integer -> Integer -> Integer
potenciaMod base 0 _ = 1
potenciaMod base exp mod
    | even exp = let p = potenciaMod base (exp `div` 2) mod
                 in (p*p) `mod` mod
    | otherwise = (base * potenciaMod base (exp - 1) mod) `mod` mod
\end{lstlisting}

\section{RSA Completo}

\begin{lstlisting}[caption={RSA en Haskell}]
generarRSA :: Integer -> Integer -> Integer -> ((Integer,Integer),(Integer,Integer))
generarRSA p q e =
    let n = p*q
        phi = (p-1)*(q-1)
        d = case inversoMod e phi of
              Just d' -> d'
              Nothing -> error "no invertible"
    in ((e,n),(d,n))

textoANumero :: String -> Integer
textoANumero = foldl (\acc c -> acc * 256 + fromIntegral (ord c)) 0

numeroATexto :: Integer -> String
numeroATexto 0 = ""
numeroATexto n = let (q,r) = n `divMod` 256
                 in numeroATexto q ++ [chr (fromIntegral r)]
\end{lstlisting}

\part{Guías de Práctica}
\chapter{Guía de Práctica 1: Técnicas Avanzadas de Conteo}

\section{Objetivos}
Al finalizar esta guía, el estudiante será capaz de:
\begin{itemize}
  \item Aplicar los principios básicos de combinatoria para resolver problemas de conteo.
  \item Establecer y resolver relaciones de recurrencia a partir de problemas del mundo real.
  \item Utilizar funciones generatrices para resolver problemas de distribución y partición.
  \item Analizar la complejidad de algoritmos recursivos e iterativos mediante relaciones de recurrencia.
\end{itemize}

\section{Principios Básicos de Combinatoria}
\subsection{Ejercicios}
\begin{enumerate}
  \item El sistema Braille para representar caracteres fue desarrollado a principios del siglo XIX por Louis Braille. Los caracteres especiales para el invidente consisten en puntos en relieve. Las posiciones para los puntos se seleccionan en dos columnas verticales de tres puntos cada una. Debe haber al menos un punto en relieve. ¿Cuántos caracteres distintos de Braille puede haber?
  \item ¿Cuántas placas de automóvil se pueden hacer que contengan tres letras seguidas de dos dígitos si se permite que haya repeticiones? ¿Y si no hay repeticiones?
  \item Encuentre el número de manos de póquer de 5 cartas (sin ordenar), seleccionadas de una baraja común de 52 cartas, que tengan las propiedades indicadas:
    \begin{enumerate}[label=\alph*)]
      \item Contienen 4 de un tipo, es decir, cuatro cartas de la misma denominación.
      \item De la forma A2345 del mismo palo (escalera de color).
      \item Contienen 2 de una denominación, 2 de otra y 1 de una tercera denominación.
    \end{enumerate}
  \item Una moneda se lanza 10 veces. ¿Cuántos resultados tienen exactamente tres caras?
  \item ¿Cuántas cadenas de 8 bits contienen a lo más 3 ceros?
  \item ¿De cuántas maneras se puede elegir un comité de 4 republicanos, 3 demócratas y 2 independientes, entre un grupo de 10 republicanos, 12 demócratas y 6 independientes?
  \item La palabra clave para acceder a un servidor de Internet está formada por 4 caracteres que se pueden elegir entre 26 letras minúsculas y 10 dígitos. Calcular el número de palabras clave que se pueden formar:
    \begin{enumerate}[label=\alph*)]
      \item Sin ninguna restricción.
      \item Usando solamente letras, sin repetirlas.
      \item Usando al menos un dígito.
    \end{enumerate}
  \item Cuatro libros distintos de matemáticas, seis diferentes de física y dos diferentes de química se colocan en un estante. ¿De cuántas formas distintas es posible ordenarlos si:
    \begin{enumerate}[label=\alph*)]
      \item Los libros de cada asignatura deben estar todos juntos.
      \item Solamente los libros de matemáticas deben estar juntos.
    \end{enumerate}
  \item Determine cuántas cadenas se pueden formar ordenando las letras ABCDEF sujetas a las condiciones indicadas:
    \begin{enumerate}[label=\alph*)]
      \item A aparece antes que D.
      \item Contiene ya sea la subcadena DB o la subcadena BE.
    \end{enumerate}
  \item ¿De cuántas formas se pueden repartir 25 rosquillas iguales entre Alan, Octavio, Carlos y Ronald si cada uno recibe al menos tres, pero no más de 8 rosquillas?
\end{enumerate}

\section{Relaciones de Recurrencia}
\subsection{Ejercicios}
\begin{enumerate}
  \item Sea $R_n$ el número de regiones en las que se divide el plano por $n$ líneas. Suponga que cada par de líneas se interseca en un punto, pero que nunca se intersecan tres líneas en un punto.
    \begin{enumerate}[label=\alph*)]
      \item Calcular $R_1, R_2, R_3$.
      \item Establezca una relación de recurrencia para la sucesión $\{R_n\}$.
    \end{enumerate}
  \item Si $\alpha$ es una cadena de bits, sea $C(\alpha)$ el número máximo de ceros consecutivos en $\alpha$. Ejemplos: $C(10010)=2$, $C(00110001)=3$. Sea $S_n$ el número de cadenas $\alpha$ de $n$ bits con $C(\alpha) \leq 2$. Desarrolle una relación de recurrencia para $S_1, S_2, \ldots$
  \item Resuelva las siguientes relaciones de recurrencia:
    \begin{enumerate}[label=\alph*)]
      \item $2a_n = 7a_{n-1} - 3a_{n-2} + 2^n$
      \item $a_n = a_{n-1} + 1 + 2^{n-1}$, con $a_0 = 2$
      \item $a_n = a_{n-1} + 6a_{n-2} + 3^n$
      \item $a_n = a_{n-1} + 1 + 2^{n-1}$, con $a_0 = 1$
      \item $a_n = 2na_{n-1}$, con $a_0 = 1$
    \end{enumerate}
  \item Resuelva la relación de recurrencia $a_n = 2n a_{n-1}$, $n>0$, con la condición inicial $a_0 = 1$.
  \item Si $\alpha$ es una cadena de $n$ bits, sea $C(\alpha)$ el número máximo de ceros consecutivos en $\alpha$ (por ejemplo, $C(0011000011) = 4$). Sea $S_n$ el número de cadenas $\alpha$ de $n$ bits con $C(\alpha) \leq 2$. Encuentre la relación de recurrencia para la sucesión $S_1, S_2, \ldots$
\end{enumerate}

\section{Funciones Generatrices}
\subsection{Ejercicios}
\begin{enumerate}
  \item ¿De cuántas formas se pueden asignar dos docenas de robots idénticos a cuatro líneas de ensamble de modo que:
    \begin{enumerate}[label=\alph*)]
      \item Al menos tres robots se asignen a cada línea.
      \item Al menos tres, pero no más de nueve robots se asignen a cada línea.
    \end{enumerate}
  \item ¿De cuántas formas pueden separarse 300 sobres idénticos en paquetes de 25 entre cuatro grupos de estudiantes, de modo que cada grupo tenga al menos 150, pero no más de 1000 sobres?
  \item Encuentre la función generatriz para el número de formas de distribuir $n$ objetos idénticos en $k$ cajas distintas, donde cada caja debe tener al menos un objeto.
  \item Utilice funciones generatrices para resolver el problema: ¿De cuántas formas se pueden elegir 5 frutas de una canasta que contiene manzanas, naranjas y peras, si hay al menos 5 de cada una?
\end{enumerate}

\section{Análisis de Algoritmos}
\subsection{Ejercicios}
\begin{enumerate}
  \item Encuentre una notación $\Theta$ para las siguientes expresiones:
    \begin{enumerate}[label=\alph*)]
      \item $(n+1)(n+3)$
      \item $2\log n + 4n + 5n\log n$
    \end{enumerate}
  \item Determine el orden para la función que cuenta el número de veces que se ejecuta la asignación en los siguientes algoritmos:
    \begin{enumerate}[label=\alph*)]
      \item \texttt{for i = 1 to n \\ for j = 1 to i \\ for k = 1 to j \\ x = x + 1}
      \item \texttt{for i = 1 to 2n \\ for j = 1 to n \\ x = x + 1}
    \end{enumerate}
  \item Dado el algoritmo recursivo para encontrar el máximo y mínimo de una sucesión (Algoritmo 10 del PDF):
    \begin{enumerate}[label=\alph*)]
      \item Sea $b_n$ el número de veces que se hace la comparación en las líneas 10 y 14. Encuentre $b_1, b_2, b_3, b_4$.
      \item Encuentre una relación de recurrencia para $(b_n)$.
      \item Resuelva la relación.
      \item Conjeture el orden de $b_n$.
    \end{enumerate}
  \item Dado el algoritmo de ordenamiento por inserción (Algoritmo 11 del PDF):
    \begin{enumerate}[label=\alph*)]
      \item Sea $b_n$ el número de veces que se hace la comparación en la línea 7 en el peor caso. Encuentre $b_1, b_2, b_3, b_4$.
      \item Encuentre una relación de recurrencia para $(b_n)$.
      \item Resuelva la relación.
      \item ¿Cuál es el orden del algoritmo?
    \end{enumerate}
  \item Algoritmo de evaluación de polinomios (Algoritmo 12 del PDF):
    \begin{enumerate}[label=\alph*)]
      \item Calcule $T(1), T(2), T(3)$.
      \item Encuentre la relación de recurrencia y condición inicial para $T(n)$.
      \item Resuelva la relación.
    \end{enumerate}
  \item Algoritmo de exponenciación rápida (Algoritmo 13 del PDF):
    \begin{enumerate}[label=\alph*)]
      \item Sea $b_n$ el número de multiplicaciones. Encuentre $b_1, b_2, b_3, b_4$.
      \item Encuentre una relación de recurrencia para $(b_n)$.
      \item Resuelva para el caso en que $n$ es potencia de 2.
      \item Conjeture la solución general.
      \item Muestre con un ejemplo que $T(n)$ es no decreciente.
      \item Demuestre que $T(n) = \Theta(\log n)$.
    \end{enumerate}
\end{enumerate}

\section{Problemas Integradores}
\begin{enumerate}
  \item Demuestre que el número de formas de distribuir $n$ objetos idénticos en $k$ cajas distintas es $\binom{n+k-1}{k-1}$.
  \item ¿Cuántos números de 5 dígitos tienen exactamente dos dígitos pares?
  \item Sea $R_n$ el número de regiones en las que se divide el plano por $n$ líneas. Suponga que cada par de líneas se interseca en un punto, pero que nunca se intersecan tres líneas en un punto.
    \begin{enumerate}[label=\alph*)]
      \item Calcular $R_1, R_2, R_3$.
      \item Establezca una relación de recurrencia para $\{R_n\}$.
    \end{enumerate}
  \item El número de formas de distribuir $n$ objetos idénticos en $k$ cajas distintas es $\binom{n+k-1}{k-1}$. Demuestre esta fórmula utilizando funciones generatrices.
  \item Analice la complejidad del algoritmo de ordenamiento por inserción y demuestre que su complejidad en el peor caso es $O(n^2)$.
  \item Demuestre que el algoritmo de exponenciación rápida (exp2) tiene complejidad $\Theta(\log n)$.
\end{enumerate}

\section{Referencias}
\begin{itemize}
  \item Richard Johnsonbaugh, \emph{Matemáticas Discretas}, 6ª ed. (2005) - Pearson Prentice Hall.
  \item Kenneth H. Rosen, \emph{Discrete Mathematics and Its Applications}, 7ª ed. (2012) - McGraw-Hill.
\end{itemize}

\section{Indicaciones para la Entrega}
\begin{itemize}
  \item La guía debe ser entregada el día del examen parcial 1.
  \item Mostrar todos los procedimientos de manera clara y ordenada.
  \item Incluir el código Haskell de las implementaciones solicitadas.
  \item Los ejercicios marcados con (*) son obligatorios para la calificación.
\end{itemize}
\chapter{Práctica: Grafos en Haskell}

\section{Objetivos}
Al finalizar esta práctica, el estudiante será capaz de:
\begin{itemize}
  \item Implementar representaciones de grafos en Haskell.
  \item Implementar algoritmos de recorrido y búsqueda.
  \item Implementar algoritmos de caminos mínimos.
\end{itemize}

\section{Ejercicios de Implementación}

\subsection{Representación y Recorridos}
\begin{enumerate}
  \item Implementar una función que convierta una matriz de adyacencia a una lista de adyacencia.
  \item Implementar DFS y BFS usando colas y pilas explícitas.
  \item Implementar una función que verifique si un grafo es conexo.
  \item Implementar una función que encuentre todas las componentes conexas.
\end{enumerate}

\subsection{Algoritmos Avanzados}
\begin{enumerate}
  \item Implementar el algoritmo de Dijkstra usando una cola de prioridad.
  \item Implementar el algoritmo de Prim.
  \item Implementar el algoritmo de Kruskal usando Union-Find.
  \item Implementar una función que encuentre un ciclo euleriano.
\end{enumerate}

\section{Proyecto}
Desarrollar una librería en Haskell para teoría de grafos con las siguientes funcionalidades:
\begin{itemize}
  \item Representación de grafos (no dirigidos, dirigidos, ponderados).
  \item Algoritmos de recorrido (DFS, BFS).
  \item Algoritmos de caminos mínimos (Dijkstra, Floyd-Warshall).
  \item Algoritmos de árboles generadores (Prim, Kruskal).
  \item Verificación de propiedades (conexidad, bipartismo, planaridad).
\end{itemize}
\chapter{Práctica: Criptografía en Haskell}

\section{Objetivos}
\begin{itemize}
  \item Implementar cifrados clásicos en Haskell.
  \item Implementar aritmética modular y criptografía RSA.
  \item Analizar la seguridad de diferentes sistemas criptográficos.
\end{itemize}

\section{Ejercicios de Implementación}

\subsection{Cifrados Clásicos}
\begin{lstlisting}[caption={Cifrados clásicos}, label=lst:ej-cripto1]
-- 1. Implementar el cifrado César con diferentes alfabetos
-- 2. Implementar el cifrado de Vigenère
-- 3. Implementar un ataque de análisis de frecuencia
-- 4. Implementar el cifrado de transposición de columnas
\end{lstlisting}

\subsection{RSA}
\begin{lstlisting}[caption={RSA}, label=lst:ej-cripto2]
-- 1. Implementar la generación de claves RSA
-- 2. Implementar cifrado y descifrado RSA
-- 3. Implementar conversión texto-número para RSA
-- 4. Implementar un ataque de factorización para claves pequeñas
\end{lstlisting}

\section{Proyecto}
Desarrollar un sistema de criptografía completo en Haskell que incluya:
\begin{itemize}
  \item Módulo de aritmética modular.
  \item Cifrados clásicos (César, Vigenère).
  \item Cifrado RSA con generación de claves.
  \item Interfaz de línea de comandos para uso práctico.
\end{itemize}

\end{document}